\documentclass[10pt,a4paper]{article}
\usepackage[utf8]{inputenc}
\usepackage{babel}
\usepackage{amsmath}
\usepackage{amsfonts}
\usepackage{amssymb}
\usepackage{amsthm}
\usepackage{mathtools}
\usepackage{microtype}
\usepackage{authblk}
\usepackage{titlesec}

\usepackage{xcolor}
\usepackage{etoolbox}
\usepackage[most]{tcolorbox}

\usepackage{empheq}
\usepackage{tikz-cd} 

% section title formating
\renewcommand{\thesection}{\Roman{section}}
\titleformat{\section}[block]{\normalsize\bfseries}{\thesection.}{1em}{\MakeUppercase}
\titlespacing*{\section}{0pt}{11pt plus 2pt minus 2pt}{11pt plus 2pt minus 2pt}

\renewcommand{\thesubsection}{\arabic{section}.\Alph{subsection}}
\titleformat{\subsection}{\normalsize\bfseries}{\thesubsection}{1em}{}
\titlespacing*{\subsection}{0pt}{11pt plus 2pt minus 2pt}{11pt plus 2pt minus 2pt}

\renewcommand{\thesubsubsection}{\thesubsection.\arabic{subsubsection}}
\titleformat{\subsubsection}{\normalsize\itshape}{\thesubsubsection}{1em}{}
\titlespacing*{\subsubsection}{0pt}{11pt plus 2pt minus 2pt}{11pt plus 2pt minus 2pt}

% text boxing
\definecolor{yaleblue}{HTML}{00356B}
\definecolor{myyellow}{HTML}{FFEEAA}

\newcommand*\resultbox[1]{%
  \colorbox{yaleblue!10}{\hspace{1em}#1\hspace{1em}}}
\newcommand*\definitionbox[1]{%
  \colorbox{myyellow!40}{\hspace{1em}#1\hspace{1em}}}

% theorem styles
\newtheoremstyle{mytheorem}
  {\topsep}   % Abstand oben
  {\topsep}   % Abstand unten
  {\itshape}  % Schriftart im Körper (kursiv)
  {}          % Einzug
  {\bfseries} % Schriftart der Überschrift (fett)
  {}          % Zeichen nach der Überschrift -> LEER (kein Punkt!)
  { }         % Abstand nach der Überschrift (ein Leerzeichen)
  {\thmname{#1}\thmnumber{ #2}\thmnote{\bfseries\itshape(#3)}}

\theoremstyle{mytheorem}
\newtheorem{theorem}{Theorem}[section]
\newtheorem{proposition}[theorem]{Proposition}
\newtheorem{corollary}[theorem]{Corollary}
\newtheorem{lemma}[theorem]{Lemma}

\newtheoremstyle{mydefinition}
  {\topsep}   % Abstand oben
  {\topsep}   % Abstand unten
  {}  % Schriftart im Körper (kursiv)
  {}          % Einzug
  {\bfseries} % Schriftart der Überschrift (fett)
  {}          % Zeichen nach der Überschrift -> LEER (kein Punkt!)
  { }         % Abstand nach der Überschrift (ein Leerzeichen)
  {\thmname{#1}\thmnumber{ #2}\thmnote{\bfseries\itshape(#3)}}

\theoremstyle{mydefinition}
\newtheorem{definition}[theorem]{Definition}

\definecolor{examplecolor}{HTML}{E0F8F8}
\newtheorem{example}[theorem]{Example}
\tcolorboxenvironment{example}{
  colback=examplecolor,
  colframe=white,
  boxrule=0pt,
  arc=0pt,
  left=6pt, right=6pt,
  top=4pt, bottom=4pt,
  breakable
}

\AtBeginEnvironment{proof}{\small}

%\newcommand{\mydef}[1]{\textls{#1}}
\newcommand{\mydef}[1]{\textbf{\textit{#1}}}

% custom macros
\newcommand{\nul}{\textrm{---}}
\newcommand{\trans}{^\mathsf{T}}

\newcommand{\ii}{\mathrm{i}}
\newcommand{\dt}{\mathrm{d}}
\newcommand{\Dt}{\mathrm{D}}
\newcommand{\Tt}{\mathrm{T}}
\newcommand{\ad}{\mathrm{ad}}
\newcommand{\Ad}{\mathrm{Ad}}

\newcommand{\SO}{\mathrm{SO}}
\newcommand{\fso}{\mathfrak{so}}
\newcommand{\fg}{\mathfrak{g}}

\newcommand{\cZ}{\mathcal{Z}}
\newcommand{\cS}{\mathcal{S}}

\newcommand{\body}{\mathrm{b}}
\newcommand{\nav}{\mathrm{n}}
\newcommand{\inert}{\mathrm{i}}

% %%%%%%%
% title %
% %%%%%%%

\title{On the analysis and connections of Lie groups}
\author{Hendrik Meier}

\affil{Oberuhldingen}

\affil{}
\renewcommand\Affilfont{\itshape}

\date{\today}


% %%%%%%%%%%
% document %
% %%%%%%%%%%

\begin{document}
\maketitle

This paper provides an introductory discussion of aspects of analysis and connection theory on Lie groups. 
The original motivation was to clarify, within the context of navigation, which connection geometrically adequately describes the transport of error vectors and covariances for attitude. 
Compared to this relatively modest motivation, the paper has grown quite lengthy---both because it aims to leave no gaps and because it takes a few detours along the way.

The reader is assumed to be familiar with the concept of a differentiable manifold. 
Furthermore, a basic familiarity with foundational concepts of differential geometry is expected, at a level typically acquired by the end of undergraduate mathematics studies.

The author makes a concerted effort to keep the notation as transparent as possible. For instance, the differential-geometric literature frequently uses the same $[-,-]$ notation for both derivatives and the algebraic Lie bracket, a practice avoided here. 
It is also the intention to clearly distinguish notations of spaces that can be canonically identified yet are not identical, and to consistently carry along the corresponding transformations explicitly.

\section{Differential Calculus on Lie Groups}

\subsection{Lie Groups}

In this paper, we treat Lie groups in the sense of the following definition.

\begin{definition}
A (real) \mydef{Lie group} $G$ is simultaneously a real $C^\infty$-manifold and a group such that group multiplication
\begin{align}
	\mathrm{mult}:
	G\times G&\rightarrow G
	\ ,\nonumber\\
	(g, h) &\mapsto gh
\end{align}
and inversion
\begin{align}
	\mathrm{inv}:
	G&\rightarrow G
	\ ,\nonumber\\
	g &\mapsto g^{-1}
\end{align}
are smooth (i.e., $C^\infty$) maps.
For multiplication, $G\times G$ inherits its $C^\infty$-structure from that of $G$.
\end{definition}

\begin{example}
\label{ex:attitude}
We describe the attitude of a rigid body in three-dimensional Euclidean space mathematically by an element $C$ of the group $\SO(3)$ of proper rotation matrices.

In the \mydef{space-fixed reference frame} $\nav$, let $e_x^\nav, e_y^\nav, e_z^\nav$ be a distinguished (right-handed) orthonormal basis; in the \mydef{body-fixed reference frame} $\body$, let this be $e_x^\body, e_y^\body, e_z^\body$.

Now let $v$ be a vector attached to a point in our ambient physical space, for example the value of a magnetic field at this location.
In space-fixed coordinates, we express $v$ by a coordinate triple
$v_\nav=(v_\nav^x, v_\nav^y, v_\nav^z)\in\mathbb{R}^3$ and in body-fixed coordinates by $v_\body=(v_\body^x, v_\body^y, v_\body^z)\in\mathbb{R}^3$, so that
\begin{align*}
	v &=
	v_\nav^x e_x^\nav
	+ v_\nav^y e_y^\nav
	+ v_\nav^z e_z^\nav
	=
	v_\body^x e_x^\body
	+ v_\body^y e_y^\body
	+ v_\body^z e_z^\body
	\ .
\end{align*}
We define the \mydef{attitude} of the rigid body with respect to the space-fixed frame as the transformation matrix $C_{\nav\body}\in\SO(3)$ that converts body-fixed coordinates into space-fixed coordinates,
\begin{align*}
	v_\nav &= C_{\nav\body} \cdot v_\body
	\ .
\end{align*}
The group $\SO(3)$, which is a Lie group by the following theorem, will accompany us repeatedly in examples throughout this paper.
\end{example}

\begin{proposition}
\label{th:SO3}
The matrix group $\SO(n)$ is a Lie group for $n\in\mathbb{N}^*$.
\end{proposition}
\begin{proof}
$\SO(n)$ is the group of $n\times n$ matrices $C$ with $CC\trans=I$ and $\det C=1$, where $I$ denotes the identity matrix.

We show that $\SO(n)$ is a submanifold of an open subset of the vector space $\mathbf{R}^{n\times n}$ of $n\times n$ matrices.
To this end, let $\mathrm{GL}_+(n)$ be the open subset in the space of $n\times n$ matrices consisting of matrices with positive determinant.
On $\mathrm{GL}_+(n)$, we consider the map
\begin{align*}
	\varphi: 
	\mathrm{GL}_+(n) &\rightarrow \mathrm{Sym}(n)
	\ ,\nonumber\\
	A &\mapsto AA\trans - I
\end{align*}
into the vector space $\mathrm{Sym}(n)$ of symmetric $n\times n$ matrices.
Its derivative is given by $\Dt\varphi|_C: A\mapsto AC\trans + CA\trans$.
For $C\in\SO(n)$, every $S\in\mathrm{Sym}(n)$ can be written as an image under $\Dt\varphi|_C$, for example as the image of $A=SC/2$, so that $\Dt\varphi|_C$ has full rank for $C\in\SO(n)$.
Hence, $\varphi$ is locally a submersion around $\SO(n)$, and $\SO(n)=\varphi^{-1}(0)$, being its fiber, is a submanifold of $\mathrm{GL}_+(n)$.

Since matrix multiplication and inversion are rational maps, both operations are smooth, and thus $\SO(n)$ is a Lie group.
\end{proof}

\subsection{Parallelization of Lie Groups}

Let $G$ be a Lie group.
Among manifolds, Lie groups are distinguished by their parallelizability.
To exploit this, we make use of right multiplication.

\begin{definition}
For $g\in G$, we call the maps
\begin{align*}
	R_g: G &\rightarrow G
	\ ,\nonumber\\
	h &\mapsto hg
\end{align*}
and
\begin{align*}
	L_g: G &\rightarrow G
	\ ,\nonumber\\
	h &\mapsto gh
\end{align*}
\mydef{right} and \mydef{left multiplication} by $g$, respectively.
\end{definition}

Using $R_g$, we parallelize $G$
by identifying tangent vectors $\xi_g\in\Tt_gG$ attached at a point $g\in G$ with a vector $\xi=\xi_e\in\Tt_eG$ in the tangent space at the identity element~$e$ according to\footnote{
Following the textbook \cite{StorchWiebe2005}, we use the symbol $\Tt$ for (coordinate-free) derivatives on manifolds.
We restrict the symbol $\Dt$ to derivatives of maps into vector spaces, in which vectors are simply differentiated component-wise after choosing a basis.
In particular, this includes $\mathfrak{g}=\Tt_eG$.
}
\begin{align}
	\xi = \Tt R_g|_e^{-1}\cdot\xi_g
	\ .
\end{align}
In this way, practical calculations with tangent vectors will take place exclusively in $\Tt_eG$,
while in formulas and derivations we still wish to distinguish transparently between $\Tt_eG$ and a tangent space $\Tt_gG$ represented by $\Tt_eG$.

\begin{definition}
\label{def:Psi1}
We denote the tangent space at the identity element by
\begin{align*}
	\mathfrak{g} := \Tt_eG
	\ .
\end{align*}
On $G\times\fg$, we define the map
\begin{align*}
	\Psi_1: G\times\fg 
	&\rightarrow\Tt G
	\ ,\nonumber\\
	(g, \xi)
	&\mapsto \Tt R_g|_e\cdot\xi
	\in \Tt_gG
	\ .
\end{align*}
\end{definition}

The map $\Psi_1$ is evidently a diffeomorphism,
since along with $R_g$, $\Tt R_g$ is also invertible (and smooth).
By means of $\Psi_1^{-1}$, $\Tt G$ is mapped diffeomorphically onto the trivial bundle $G\times\fg$.

\begin{example}
Let $\varphi:A\mapsto AA\trans - I$ be the submersion $\mathrm{GL}_+(3)\rightarrow\mathrm{Sym}(3)$ already used in the proof of \ref{th:SO3}, whose fiber $\varphi^{-1}(0)$ yields $\SO(3)$ as a submanifold.

We call the tangent space $\Tt_I\SO(3)$ at the identity matrix, the identity element of $\SO(3)$,
\begin{align*}
	\fso(3) := \Tt_I\SO(3)
	\ .
\end{align*}
We obtain it as a subspace of $\Tt_I\mathrm{GL}_+(3)=\mathbb{R}^{3\times 3}$ in the form of the kernel of $\Dt\varphi|_I:A\mapsto A+A\trans$ and find
\begin{align*}
	\fso(3) &= 
	\big\{
		A\in\mathbb{R}^{3\times 3}
		\big| 	
		A = -A\trans
	\big\}
	\ .
\end{align*}

Let $C\in\SO(3)$.
In matrix space, the derivative of right multiplication $R_C:Q\mapsto QC$ is simply computed as $\Tt R_C|_Q:A_Q\mapsto A_QC$ for $A_Q\in\Tt_Q\SO(3)$.
In particular, it follows that
\begin{align*}
	\Tt_C\SO(3) &= \fso(3)\cdot C
	\ ,
\end{align*}
where $\Tt_C\SO(3)$ is also regarded as a subspace of $\Tt_C\mathrm{GL}_+(3)$ ($\cong\mathbb{R}^{3\times 3}$).

Accordingly, for $\SO(3)$, the map $\Psi_1$ takes the form
\begin{align*}
	\Psi_1(C, A) &= AC
	\ .
\end{align*}
\end{example}

Let $\xi\in\mathfrak{g}=\Tt_eG$ be a tangent vector attached at the identity element.
Via $\Psi_1$, we define the vector field
\begin{align}
	X_\xi: G &\rightarrow \Tt G
	\ ,\nonumber\\
	g &\mapsto \Psi_1(g, \xi)\in\Tt_g G
	\ .
\end{align}
It is right-invariant in the sense of the following definition.

\begin{definition}
We call a vector field $X$ on $G$ \mydef{right-invariant} if and only if
\begin{align}
	X|_{hg} &= \Tt R_g|_h\cdot X|_{h}
\end{align}
holds for all $g\in G$.
\end{definition}

A right-invariant vector field~$X$ is uniquely determined by its value $X|_e$ at the identity element and can therefore be identified with the tangent vector $\xi=X|_e\in\mathfrak{g}$.

\begin{example}
\label{ex:attitude_error}
Let the attitude of a body-fixed frame in the space-fixed frame be described by $C\in\SO(3)$ as in Example~\ref{ex:attitude}.

We interpret a tangent vector $A_C$ in $\Tt_C\SO(3)$ as a \mydef{virtual displacement}, which describes, for instance, a (possible) error in our attitude given by $C$.
(In applications, this is the case when $C$ merely represents an \emph{estimate}.)

Now suppose the attitude is subject to motion.
Differentiating $CC\trans=I$ yields $\Omega(t)=C(t)\trans\dot{C}(t)\in\fso(3)$.
$\Omega(t)$ is the body-fixed angular velocity, measurable for example by a gyroscope. (The assumption is that the space-fixed frame does not rotate.)

If $\Omega(t)$ is the ideal---that is, error-free and as a continuous function of time---measured angular velocity,
we can track the change in attitude without error by solving the differential equation
\begin{align*}
	\dot{\tilde{C}}(t)
	&= \tilde{C}(t)\cdot\Omega(t)
\end{align*}
with initial condition $\tilde{C}(0)=I$.
If the initial attitude is given by $C(0)$,
then $C(t)=C(0)\tilde{C}(t)$.
Thus, the propagation of attitude proceeds via \emph{right multiplication}.

If $C(0)$ is only available as an estimate, we must account for an error represented by a tangent vector $A_{C(0)}(0)\in\Tt_{C(0)}\SO(3)$.
In our \emph{interpretation}, this error is meant to indicate the extent to which the \emph{space-fixed} frame estimated relative to the body via $C(0)$ deviates from the true space-fixed (reference) frame.

Since an error residing on the space-fixed side does not change under the ideally tracked motion, we interpret it along the curve $C(t)$ as a \emph{right-invariant} vector field.
In practice, we consistently keep track of the error as $A(0)\in\fso(3)$ via
\begin{align*}
	A(0) := A_I(0) &= A_{C(t)}(0)\cdot C\trans(0)
	\ .
\end{align*}
\end{example}

\subsection{Differentiation on Lie Groups}

The derivative of a differentiable map $F:G\rightarrow G$ at a point $g\in G$ is initially the map
\begin{align*}
	\Tt F|_g:
	\Tt_gG\rightarrow\Tt_{F(g)}G
	\ .
\end{align*}
For practical calculations, we identify both the domain and image with $\mathfrak{g}=\Tt_eG$ using the parallelization realized by $\Psi_1$.
We denote the derivative treated in $\mathfrak{g}$ in this manner by ``$\Dt$''.

\begin{definition}
We define
\begin{align*}
	\Dt F|_g: 
	\mathfrak{g} &\rightarrow \mathfrak{g}
	\ ,\nonumber\\
	\xi &\mapsto 
	\Psi_1^{-1}\big(
		\Tt F|_g
		\cdot \Psi_1(g, \xi)
	\big)
	=
	\Tt R_{F(g)^{-1}}|_{F(g)}
	\cdot \Tt F|_g
	\cdot \Tt R_{g}|_e\cdot\xi
\end{align*}
to be the derivative $\Tt F|_g$ transformed via $\Psi_1$ into a map $\fg\rightarrow\fg$.\footnote{
	Evidently, the outermost factor in the composition can also be written as $\Tt R_{F(g)^{-1}}|_{F(g)}=\Tt R_{F(g)}^{-1}|_{F(g)}=\Tt R_{F(g)}|_e^{-1}$.
}
\end{definition}

\subsubsection{Product Rule}

For the derivative of Lie group multiplication
\begin{align}
	\mathrm{mult}:G\times G\rightarrow G
	\ ,
\end{align}
we distinguish between the derivative
\begin{align*}
	\Tt\,\mathrm{mult}|_{(g,h)}:\Tt_gG\times\Tt_hG\rightarrow\Tt_{hg}G
\end{align*}
and the derivative
\begin{align*}
	\Dt\,\mathrm{mult}|_{(g,h)}:\fg\times\fg\rightarrow\fg
	\ .
\end{align*}

\begin{lemma}
\label{th:Tmult}
We have
\begin{align*}
	\Tt\,\mathrm{mult}|_{(g,h)}\cdot (\xi_g, \upsilon_h)
	&= \Tt R_h|_g\cdot\xi_g + \Tt L_g|_h\cdot\upsilon_h
	\ .
\end{align*}
\end{lemma}

\begin{proof}
For $\varepsilon>0$, let $\gamma,\eta: (-\varepsilon,\varepsilon)\rightarrow G$ be curve germs with $\gamma(0)=g$, $\dot{\gamma}(0)=\xi_g\in\Tt_gG$, as well as $\eta(0)=h$, $\dot{\eta}(0)=\upsilon_h\in\Tt_hG$.

With these, we find on the one hand
\begin{align*}
	\Tt\,\mathrm{mult}|_{(g,h)}\cdot (\xi_g, 0)
	&= \frac{\dt}{\dt t}
		R_h\gamma(t)
	\Big|_{t=0}
	= \Tt R_h|_g\cdot\xi_g
	\in \Tt_{gh}G
	\ ,
\end{align*}
and on the other hand
\begin{align*}
	\Tt\,\mathrm{mult}|_{(g,h)}\cdot (0, \upsilon_h)
	&= \frac{\dt}{\dt t}
		L_g\eta(t)
	\Big|_{t=0}
	= \Tt L_g|_h\cdot\upsilon_h
	\in \Tt_{gh}G
	\ .
\end{align*}
Combining the two yields the claim.
\end{proof}

\begin{proposition}
\label{th:Dmult}
We have
\begin{align*}
	\Dt\,\mathrm{mult}|_{(g, h)} \cdot (\xi, \upsilon)
	&=
	\xi + \Tt(R_{g^{-1}}\circ L_g)|_e\cdot\upsilon
	=
	\xi + \Ad(g)\cdot\upsilon
	\ ,
\end{align*}
where $\Ad$ is given by the following definition.
\end{proposition}

\begin{definition}
\label{def:Ad}
The Lie group homomorphism
\begin{align}
	\Ad: G&\rightarrow \mathrm{GL}(\mathfrak{g})
	\ ,\nonumber\\
	g &\mapsto
	\Tt
	(R_{g^{-1}}\circ L_g)|_e
\end{align}
is called the \mydef{adjoint representation} of $G$.
\end{definition}

That $\Ad$ is a Lie group homomorphism is verified immediately.

\begin{proof}[Proof of \ref{th:Dmult}]
We pick up at the end of the proof of \ref{th:Tmult} and compute
\begin{align*}
	\Dt\,\mathrm{mult}|_{(g,h)}\cdot (\xi, 0)
	&= 
	\Tt R_{gh}^{-1}|_{gh}
	\cdot
	\Tt R_h|_g
	\cdot
	\Tt R_g|_g
	\cdot\xi
	\nonumber\\
	&= \Tt (R_{(gh)^{-1}}\circ R_h\circ R_g)|_e\cdot\xi
	= \xi
\end{align*}
as well as
\begin{align*}
	\Dt\,\mathrm{mult}|_{(g,h)}\cdot (0, \upsilon)
	&=
	\Tt R_{gh}^{-1}|_{gh}
	\cdot	
	\Tt L_g|_h
	\cdot
	\Tt R_h|_e
	\cdot
	\upsilon
	=
	\Tt
	(R_{g^{-1}}\circ L_g)|_e
	\cdot
	\upsilon	
	\ .
\end{align*}
Combining the two yields the claim.
\end{proof}

\begin{example}
\label{ex:Ad}
Let $C\in\SO(3)$ and $A\in\fso(3)$.
In the adjoint representation, $C$ acts on $A$ according to
\begin{align*}
	\Ad(C)\cdot A &= CAC\trans
	\ .
\end{align*}
Since $A$ is skew-symmetric, $CAC\trans$ is evidently skew-symmetric as well.

Let $a=(a^1, a^2, a^3)\in\mathbb{R}^3$.
Using the standard vector product, we define the isomorphism
\begin{align*}
	[-]_\times: 
	\mathbb{R}^3&\rightarrow\fso(3)
	\ ,\nonumber\\
	a &\mapsto [a]_\times = a\times-
	\ .
\end{align*}
That is,
\begin{align*}
	[a]_\times
	&=
	\begin{pmatrix}
		0    & -a^3 &  a^2 \\
		 a^3 &    0 & -a^1 \\
		-a^2 &  a^1 &    0
	\end{pmatrix}
	\in\fso(3)
	\ .
\end{align*}
From
\begin{align*}
	C[a]_\times C\trans &= [Ca]_\times
\end{align*}
it follows that in $\mathbb{R}^3$, identified with $\fso(3)$ via $[-]_\times$, the adjoint representation corresponds to the matrix-vector product $Ca$.
\end{example}

\begin{example}
\label{ex:attitude_mult}
We adopt the interpretation from Example~\ref{ex:attitude_error} of tangent vectors as virtual displacements representing errors.
In addition to the reference frames $\body$ and $\nav$, suppose there is a further frame $\inert$.
For the attitudes between $\body$ and $\nav$ as well as between $\nav$ and $\inert$, suppose we are given estimates $C_{\nav\body}$ and $C_{\inert\nav}$ that are subject to errors.
Let these errors be given by $A_{\nav\body}$ and $A_{\inert\nav}$ in $\fso(3)$, respectively.

From this information, the estimate for the attitude of $\body$ relative to $\inert$ can be written as
\begin{align*}
	C_{\inert\body}
	&= C_{\inert\nav} C_{\nav\body}
	\ .
\end{align*}
The error $A_{\inert\body}\in\fso(3)$ inherent in this product propagates from the errors of the factors according to the product rule~\ref{th:Dmult},
\begin{align*}
	A_{\inert\body}
	&=
	\Dt\,\mathrm{mult}|_{(C_{\inert\nav}, C_{\nav\body})}
	\cdot (A_{\inert\nav}, A_{\nav\body})
	\nonumber\\
	&=
	A_{\inert\nav} + C_{\inert\nav}A_{\nav\body}C_{\inert\nav}\trans
	\ .
\end{align*}
If, as in Example~\ref{ex:Ad}, we identify the space $\fso(3)$ with $\mathbb{R}^3$, we find, in notation adapted according to obvious rules, that
\begin{align*}
	a_{\inert\body}
	&=
	a_{\inert\nav} + C_{\inert\nav}a_{\nav\body}
	\ .
\end{align*}
Consequently, the error of the inner attitude estimate must first be rotated by the outer attitude estimate before being combined with the latter's error.
\end{example}

\subsubsection{Derivative of Inversion}

For the derivative of inversion, the product rule yields the following corollary.

\begin{corollary}
\label{th:DInv}
We have
\begin{align*}
	\Tt\,\mathrm{inv}|_{g}
	&= -\Tt (R_{g^{-1}}\circ L_{g^{-1}})|_g
\end{align*}
for the derivative $\Tt\,\mathrm{inv}|_{g}:\Tt_gG\rightarrow\Tt_{g^{-1}}G$, and
\begin{align*}
	\Dt\,\mathrm{inv}|_{g}
	&= -\Ad(g^{-1})
\end{align*}
for the derivative $\Dt\,\mathrm{inv}|_{g}:\mathfrak{g}\rightarrow\mathfrak{g}$.
\end{corollary}

\begin{example}
For a curve $C:\mathbb{R}\rightarrow\SO(3)$, we have
\begin{align*}
	\dot{C}\trans &= -C\trans\dot{C}C\trans
	\ ,
\end{align*}
as obtained immediately by differentiating $CC\trans=I$.

Setting $\Omega_\nav=\dot{C}C\trans\in\fso(3)$, we find for $\Omega_\nav'=\dot{C}\trans C$ that
\begin{align*}
	\Omega_\nav' &= -C\trans\Omega_\nav C = -\Ad(C\trans)\cdot \Omega_\nav
	\ .
\end{align*}
\end{example}

\begin{proof}[Proof of \ref{th:DInv}]
We express the \emph{constant} map to the identity element as the composition
\begin{alignat}{2}
	G 
	&\rightarrow G\times G 
	&&\rightarrow G
	\ ,\nonumber\\
	g 
	&\mapsto
	\begin{pmatrix}
		\mathrm{inv}(g) \\ g
	\end{pmatrix}
	&&\mapsto \mathrm{mult}(\mathrm{inv}(g), g) = e
	\ .
\end{alignat}
For $g\in G$ and $\xi_g\in\Tt_gG$, the chain rule then gives
\begin{align*}
	\mathfrak{g}\ni	
	0 &= 
	\Tt\,\mathrm{mult}|_{(g^{-1},g)}
	\cdot
	\begin{pmatrix}
		\Tt\,\mathrm{inv}|_{g} \cdot \xi_g \\
		\xi_g
	\end{pmatrix}
	\nonumber\\
	&= 
	\Tt R_g|_{g^{-1}}\cdot \Tt\,\mathrm{inv}|_{g} \cdot \xi_g
	+ \Tt L_{g^{-1}}|_g\cdot \xi_g
	\ .
\end{align*}
With $\Tt R_g|_{g^{-1}}^{-1}=\Tt R_{g^{-1}}|_{e}$, the formula for $\Tt\,\mathrm{inv}$ follows.

On the other hand, substituting $\xi_g=\Tt R_g|_e\cdot \xi$, we find
\begin{align*}
	\Dt\,\mathrm{inv}|_{g} \cdot \xi
	&= \Tt R_g|_{g^{-1}}
	\cdot \Tt\,\mathrm{inv}|_{g}
	\cdot \Tt R_g|_e\cdot \xi
	\nonumber\\
	&= -\Tt L_{g^{-1}}|_g\cdot \Tt R_g|_e \cdot \xi
	= -\Tt (L_{g^{-1}}\circ R_g)|_e \cdot \xi
	\ ,
\end{align*}
which yields the formula for $\Dt\,\mathrm{inv}$.
\end{proof}

\subsection{Vector Fields on Manifolds}

In this section, we collect basic concepts for working with vector fields on manifolds.\footnote{
	Throughout, we consider smooth ($C^\infty$) manifolds.
}
In what follows, let $X$ generally be a vector field on a manifold $M$.

\subsubsection{Push-Forward}

Via diffeomorphisms, vector fields can be transferred within a manifold or between manifolds.
This is accomplished by the push-forward in the sense of the following definition.
\begin{definition}
Let $F:M\rightarrow N$ be a diffeomorphism onto the manifold~$N$.
The \mydef{push-forward} of $X$ onto $N$ is the vector field defined on $N$ by
\begin{align}
	F_*X: N &\rightarrow \Tt N
	\ ,\nonumber\\
	q &\mapsto \Tt F|_{F^{-1}(q)}\cdot X|_{F^{-1}(q)}
	\ .
\end{align}
\end{definition}

\subsubsection{Flows}

Of central importance in the study of vector fields are their integral curves and flows.
\begin{definition}
We call a curve $t\mapsto\gamma(t)\in M$ an \mydef{integral curve} of the vector field~$X$ if its instantaneous velocity is given by the vector field at its current position,
\begin{align*}
	\dot{\gamma}(t) &= X|_{\gamma(t)}
	\ .
\end{align*}
\end{definition}

\begin{theorem}
There exist an open interval $I\subset\mathbb{R}$ with $0\in I$ and, for each $p_0\in M$, an open set $V\subset M$ with $p_0\in V$, on whose product the map
\begin{align}
	\Phi_X: I\times V \rightarrow M
\end{align}
satisfying
\begin{align}
	\frac{\dt}{\dt t}\Phi_X(t; p) = X|_{\Phi_X(t; p)}
\end{align}
and $\Phi_X(0; p)=p$ is well-defined and smooth.
\end{theorem}
\begin{proof}
The statement follows from the theory of ordinary differential equations.
\end{proof}

\begin{definition}
The map $\Phi_X$ of the preceding theorem is called the \mydef{flow} of $X$.
\end{definition}

For fixed $p$, the curve $t\mapsto\Phi_X(t; p)$ is the unique integral curve of $X$ passing through $p$.
The flow $\Phi_X$ of the (time-independent) vector field~$X$ evidently satisfies the \mydef{semigroup property} in the sense that
\begin{align}
	\Phi_X(t_1 + t_2; -)
	&= \Phi_X(t_2; \Phi_X(t_1; -))
	\ ,
\end{align}
whenever the right-hand side exists.

\subsubsection{Lie Derivative of Vector Fields}

For $t\in I$, $\Phi_X(t):=\Phi_X(t;-)$ is (at least) \emph{locally} a diffeomorphism onto its image.
We denote its (local) inverse by $\Phi_X^{-1}(t):q\mapsto\Phi_X^{-1}(t;q)$, where $\Phi_X^{-1}(t;q)=\Phi_X(-t;q)$.
Now let $Y$ be another vector field on $M$.

\begin{definition}
The vector field defined for all $p\in M$ by
\begin{align}
	\{X, Y\}\big|_{p} &:= 
	\frac{\dt}{\dt t}
	\big(
		[\Phi_X(t)]_* Y
	\big)
	\big|_{p}
	\Big|_{t=0}
\end{align}
is called the \mydef{Lie derivative} (or \mydef{Lie product}) of the vector fields $X$ and $Y$.
\end{definition}

\begin{definition}
The vector field defined in a chart of~$M$ by
\begin{align*}
	[X, Y]\big|_{p}
	&=
	\Dt X|_{p}
	\cdot
	Y|_{p}
	-
	\Dt Y|_{p}
	\cdot
	X|_{p}
\end{align*}
is called the \mydef{Lie bracket} of the vector fields $X$ and $Y$.
\end{definition}
That the Lie bracket is independent of the choice of chart follows from
\begin{lemma}
\label{th:pushForwardLie}
Let $M=\mathbb{R}^n$ and let $F:M\rightarrow\mathbb{R}^n$
be a diffeomorphism.
Then
\begin{align*}
	[F_*X, F_*Y]\big|_{p}
	&= F_*[X, Y]|_p
	\ .
\end{align*}
\end{lemma}
Moreover, we have
\begin{theorem}
\label{th:LieBracket}
The Lie derivative and the Lie bracket coincide,
\begin{align*}
	\{X, Y\} &= [X, Y]
	\ .
\end{align*}
\end{theorem}
In the following, we use the common notation~$[-,-]$ for both.
Equipped with the Lie bracket (or Lie product), the vector space of vector fields on $M$ becomes an (infinite-dimensional) Lie algebra.
The Jacobi identity for the Lie bracket is confirmed by direct verification.

\begin{proof}[Proof of \ref{th:LieBracket}]
To prove this local statement, we assume without loss of generality that $M=\mathbb{R}^n$.
For
\begin{align*}
	\{X, Y\}\big|_{p}
	&=
	\frac{\dt}{\dt t}
	\Big(
		\Dt\Phi_X(t)|_{\Phi_X^{-1}(t; p)}
		\cdot
		Y|_{\Phi_X^{-1}(t; p)}
	\Big)\Big|_{t=0}
\end{align*}
we find, using
\begin{align*}
	\Phi_X(t; p)
	&= p + t X|_{p} + \mathrm{O}(t^2)
	\ ,\nonumber\\
	\Phi_X^{-1}(t; p)
	&= p - t X|_{p} + \mathrm{O}(t^2)
	\ ,\nonumber\\
	\Dt\Phi_X(t)\big|_{p}
	&= \mathrm{id} + t\,\Dt X|_{p} + \mathrm{O}(t^2)
	\ ,\nonumber\\
	Y|_{\Phi_X^{-1}(t; p)}
	&= Y|_{p} - t\,\Dt Y|_{p}\cdot X|_{p} + \mathrm{O}(t^2)
	\ ,
\end{align*}
that
\begin{align*}
	\{X, Y\}\big|_{p}
	&=
	\Dt X|_{p}
	\cdot
	Y|_{p}
	-
	\Dt Y|_{p}
	\cdot
	X|_{p}
	\ .
\end{align*}
\end{proof}

\begin{proof}[Proof of \ref{th:pushForwardLie}]
By direct calculation, we find
\begin{align*}
	[F_*X, F_*Y]\big|_{p}
	&= 
	\Dt \big(
		\Dt F|_{F^{-1}(p)} \cdot X|_{F^{-1}(p)}
	\big)\big|_p
	\cdot
	\Dt F|_{F^{-1}(p)} \cdot Y|_{F^{-1}(p)}
	- [X\leftrightarrow Y]
	\nonumber\\
	&=
	\big[
		\Dt^2F|_{F^{-1}(p)}
		\cdot
		(X|_{F^{-1}(p)}\otimes Y|_{F^{-1}(p)})
		\nonumber\\
		&\quad 
		+
		\Dt F|_{F^{-1}(p)} 
		\cdot \Dt X|_{F^{-1}(p)}
		\cdot Y|_{F^{-1}(p)}
	\big]
	- [X\leftrightarrow Y]
	\ .
\end{align*}
Here we have used that for vectors $u,v\in\mathbb{R}^n$,
\begin{align*}
	\Dt_v(\Dt_u F|_{F^{-1}(p)})\big|_p
	&=
	\Dt_{\Dt F^{-1}|_p\cdot v}
	\Dt_u F\big|_{F^{-1}(p)}
	=:
	\Dt^2F|_{F^{-1}(p)}\cdot(u \otimes \Dt F|_{F^{-1}(p)}^{-1}\cdot v)
\end{align*}
with $\Dt F^{-1}|_p=\Dt F|_{F^{-1}(p)}^{-1}$, and that
\begin{align*}
	\Dt \big(
		X|_{F^{-1}(p)}
	\big)\big|_p
	&=
	\Dt X|_{F^{-1}(p)}
	\cdot \Dt F|_{p}^{-1}
	\ .
\end{align*}
For the antisymmetrized expression, it follows that
\begin{align*}
	[F_*X, F_*Y]\big|_{p}
	&= \Dt F|_{F^{-1}(p)}
	\cdot
	[X, Y]|_{F^{-1}(p)}
	= F_*[X, Y]|_p
\end{align*}
since $\Dt^2$ is symmetric.
\end{proof}

\subsection{Lie Algebras}

Let $G$ be a Lie group again.

\subsubsection{Flow of a Right-Invariant Vector Field}

Let $\xi\in\mathfrak{g}$ and let $X_\xi:g\mapsto\Psi_1(g,\xi)=\Tt R_g|_e\cdot\xi$ be the associated right-invariant vector field.
For $\varepsilon>0$, let $\gamma:(-\varepsilon,\varepsilon)\rightarrow G$ be a local integral curve of $X_\xi$ passing through the identity element,
\begin{align*}
	\gamma(0) &= e
	\ ,\quad
	\dot{\gamma}(t) = X_\xi|_{\gamma(t)}
	\ .
\end{align*}

\begin{proposition}
\label{th:FlowL}
In the given situation,
\begin{align*}
	\Phi_{X_\xi}(t; -)
	&= L_{\gamma(t)}
\end{align*}
is the flow of the right-invariant vector field $X_\xi$.
\end{proposition}
\begin{proof}
We have
\begin{align*}
	\frac{\dt}{\dt t}
	\big(
		L_{\gamma(t)}\,g
	\big)
	\Big|_{t=0}
	&=
	\frac{\dt}{\dt t}
	\big(
		\gamma(t)\,g
	\big)
	\Big|_{t=0}	
	=
	\Tt R_g|_{e} \cdot X_\xi|_{e}
	= X_\xi|_{g}		
\end{align*}
and $L_{\gamma(0)}=L_{e}=\mathrm{id}$.
\end{proof}

\begin{example}
\label{ex:flowSOn}
On $\SO(n)$ for $n\geq 2$, consider the curve $C(t)=\exp_\mathsf{M}(tA)$ with $A\in\fso(n)$.
Here, $\exp_\mathsf{M}$ denotes the matrix exponential defined by the power series of $x\mapsto\mathrm{e}^x$.
Evidently, $C(0)=I$ and $\dot{C}(t)=A\exp_\mathsf{M}(tA)=AC(t)$. The expression $AC(t)$ equals the value of the right-invariant vector field~$X_A$ defined by $A$ at the point $C(t)$.

Thus, the flow of $X_A$ with $X_A|_I=A$ has the form
\begin{align*}
	\Phi_{X_A}(t; Q) &= \exp_\mathsf{M}(tA)\,Q
	\ .
\end{align*}
\end{example}

\subsubsection{Lie Bracket on $\mathfrak{g}$}

Using the push-forward, right-invariant vector fields can also be characterized as follows.
\begin{lemma}
A vector field $X$ is right-invariant if and only if
\begin{align*}
	(R_g)_*X &= X
\end{align*}
holds for all $g\in G$.
\end{lemma}
\begin{proof}
Let $h\in G$. We have $[(R_g)_*X]|_h=\Tt R_g|_{hg^{-1}}\cdot X|_{hg^{-1}}=X|_h$, where the second equality corresponds to right-invariance.
\end{proof}

Let $\xi_1,\xi_2\in\mathfrak{g}=\Tt_eG$ be two tangent vectors at the identity element with associated right-invariant vector fields~$X_{\xi_1}$ and $X_{\xi_2}$.

\begin{proposition}
The Lie bracket $[X_{\xi_1}, X_{\xi_2}]$ is a right-invariant vector field.
\end{proposition}
\begin{proof}
From \ref{th:pushForwardLie} and the preceding lemma, it follows that
\begin{align*}
	(R_g)_*[X_{\xi_1}, X_{\xi_2}]
	&= [(R_g)_*X_{\xi_1}, (R_g)_*X_{\xi_2}]
	= [X_{\xi_1}, X_{\xi_2}]
	\ .
\end{align*}
Applying this lemma once more yields the claim.
\end{proof}
Consequently, the Lie bracket of two right-invariant vector fields can again be identified with a tangent vector in $\fg=\Tt_eG$.
We have thus shown:
\begin{theorem}
The tangent space $\fg=\Tt_eG$ at the identity element of a Lie group $G$ carries the structure of a \mydef{Lie algebra}.
The Lie bracket on $\fg$ is given by
\begin{align*}
	[-,-]^\mathfrak{g}:
	\mathfrak{g}\times\mathfrak{g}
	&\rightarrow \mathfrak{g}
	\ ,\nonumber\\
	(\xi_1,\xi_2) &\mapsto
	[\xi_1,\xi_2]^\mathfrak{g}
	:= [X_{\xi_1}, X_{\xi_2}]|_e
	\ .
\end{align*}
\end{theorem}
$[-,-]^\fg$ is a bilinear map on $\mathfrak{g}$ taking values in $\mathfrak{g}$.
The superscript $\fg$ is intended to distinguish the algebraic Lie bracket on $\fg$ from the differential Lie bracket $[-,-]$ between arbitrary vector fields.

\begin{example}
Let $A_1, A_2\in\fso(n)$ with associated right-invariant vector fields $X_{A_i}|_C=A_iC$ for $i=1,2$.
We compute their Lie bracket according to
\begin{align*}
	[A_1, A_2]^{\fso(n)}
	&= \frac{\dt}{\dt t}
		\big(
			[\Phi_{X_{A_1}}(t)]_* X_{A_2}
		\big)
		\big|_{I}
	\Big|_{t=0}
	\ .
\end{align*}
Substituting the flow from Example~\ref{ex:flowSOn} yields
\begin{align*}
	[\Phi_{X_{A_1}}(t)]_* X_{A_2}
	&= \exp_\mathsf{M}(tA_1)\,A_2\,\exp_\mathsf{M}(-tA_1)
	\ ,
\end{align*}
from which differentiating at $t=0$ gives
\begin{align*}
	[A_1, A_2]^{\fso(n)} = A_1A_2-A_2A_1
	\ .
\end{align*}
The Lie bracket on $\fso(n)$ is therefore equal to the matrix commutator.
\end{example}

Specifically, let $\xi_1,\ldots,\xi_n$ be a basis of $\mathfrak{g}$.
Then there is a representation
\begin{align}
	[\xi_i,\xi_j]^\mathfrak{g}
	&= \sum_{k=1}^n f^k_{ij} \xi_k
	\ .
\end{align}
The \mydef{structure constants} $f^k_{ij}\in\mathbb{R}$ with respect to a given basis uniquely determine the Lie algebra structure of $\mathfrak{g}$.

\begin{example}
\label{ex:so3base}
For $\fso(3)$, the vector space of skew-symmetric $3\times 3$ matrices, we choose the basis
\begin{align*}
A_1 &=
\begin{pmatrix}
0 & 0 & 0 \\
0 & 0 & -1 \\
0 & 1 & 0
\end{pmatrix}
\ ,\quad
A_2 =
\begin{pmatrix}
0 & 0 & 1 \\
0 & 0 & 0 \\
-1& 0 & 0
\end{pmatrix}
\ ,\quad
A_3 =
\begin{pmatrix}
0 & -1 & 0 \\
1 & 0 & 0 \\
0 & 0 & 0
\end{pmatrix}
\ .
\end{align*}
For this basis, we have
\begin{align*}
[A_i,A_j]^{\fso(3)} &= \sum_{k=1}^3 \varepsilon_{ijk} A_k
\ ,
\end{align*}
so that the structure constants $f^k_{ij}=\varepsilon_{ijk}$ with respect to the given basis coincide with the Levi-Civita symbol.

By inversion, the vector space isomorphism $[-]_\times:\mathbb{R}^3\rightarrow\fso(3)$ assigns to the basis vector $A_i$ the $i$-th vector $e_i$ of the standard basis of $\mathbb{R}^3$.
At the same time, the standard vector product satisfies the relations $e_i\times e_j=\sum_{k=1}^3 \varepsilon_{ijk}e_k$.

As a result, the isomorphism $a\mapsto[a]_\times$ is even a Lie algebra isomorphism $(\mathbb{R}^3,\times)\rightarrow(\fso(3), [-,-]^{\fso(3)})$.
\end{example}

\subsubsection{Adjoint Representation}

The adjoint representation $\Ad:G\rightarrow\mathrm{GL}(\mathfrak{g})$ according to Definition~\ref{def:Ad} is a Lie group homomorphism.
Its derivative evaluated at the identity element plays a similar role at the level of the Lie algebra.

\begin{theorem}
\label{th:ad}
The derivative of the adjoint representation evaluated at the identity element $e$,
\begin{align*}
	\ad:=\Dt\,\Ad|_e: \mathfrak{g}\rightarrow\mathfrak{gl}(\mathfrak{g})
	\ ,
\end{align*}
is the derivation on~$\mathfrak{g}$ given by
\begin{align*}
	\ad(\xi): \upsilon\mapsto[\xi,\upsilon]^\fg
	\ .
\end{align*}
\end{theorem}
That $\ad(\xi)$ is a derivation follows from the Jacobi identity on $\mathfrak{g}$ and $\upsilon_1,\upsilon_2\in\fg$ via
\begin{align*}
\ad(\xi) [\upsilon_1,\upsilon_2]^\fg
&= 
[[\xi, \upsilon_1]^\fg,\upsilon_2]^\fg + [\upsilon_1, [\xi, \upsilon_2]^\fg]^\fg
=[\ad(\xi) \upsilon_1,\upsilon_2]^\fg + [\upsilon_1,\ad(\xi)\upsilon_2]^\fg
\ .
\end{align*}
Likewise, from the Jacobi identity, it follows via
\begin{align*}
	\ad([\xi_1,\xi_2]^\fg)
	&= [[\xi_1,\xi_2]^\fg, -]^\fg
	= -[[\xi_2,-]^\fg, \xi_1]^\fg-[[-,\xi_1]^\fg, \xi_2]^\fg
	\nonumber\\
	&= \ad(\xi_1)\circ \ad(\xi_2) - \ad(\xi_2)\circ \ad(\xi_1)
	\ ,
\end{align*}
that $\ad$ is a Lie algebra homomorphism $\mathfrak{g}\rightarrow\mathfrak{gl}(\mathfrak{g})$.

\begin{proof}[Proof of \ref{th:ad}]
By definition, for $g\in G$ and $\upsilon\in\mathfrak{g}$ with associated right-invariant vector field~$Y_\upsilon$, we have
\begin{align*}
	\Ad(g)\cdot\upsilon
	&= 
	\Tt(L_{g}R_{g^{-1}})|_e\,\upsilon
	=
	\Tt L_{g}|_{g^{-1}}
	\cdot
	\upsilon_{g^{-1}}
	=  
	\big(
		[L_{g}]_* Y_\upsilon
	\big)
	\big|_e
	\ .	
\end{align*}
Let $\gamma:(-\varepsilon,\varepsilon)\rightarrow G$ with $\gamma(0)=e$ and $\dot{\gamma}(t)=X_\xi|_{\gamma(t)}$ be an integral curve of the right-invariant vector field $X_\xi$ with $X_\xi|_e=\xi$.
It follows that
\begin{align*}
	\Dt_{\xi}\Ad|_e\cdot\upsilon
	&=
	\frac{\dt}{\dt t}
	\Ad(\gamma(t))\cdot\upsilon
	\Big|_{t=0}
	=
	\frac{\dt}{\dt t}
	\big(
		[L_{\gamma(t)}]_* Y_\upsilon
	\big)
	\big|_e
	\Big|_{t=0}
	\ .
\end{align*}
With $L_{\gamma(t)}=\Phi_{X_\xi}(t;-)$ by Lemma~\ref{th:FlowL}, it follows that
\begin{align*}
	\Dt_{\xi}\Ad|_e\cdot\upsilon
	&=
	[X_\xi,Y_\upsilon]|_e = [\xi, \upsilon]^\mathfrak{g}
	\ .
\end{align*}
\end{proof}
	
\subsection{Lie Exponential Map}

Let $\xi\in\mathfrak{g}=\Tt_eG$ be an element of the Lie algebra of the Lie group~$G$ and $X_\xi$ the associated right-invariant vector field.
From the theory of ordinary differential equations, it follows that the flow $\Phi_{X_\xi}(t;g)$ is defined at least locally for sufficiently small $t$.
This statement can be strengthened:
\begin{theorem}
\label{th:PhiGlobal}
The flow $(t, g)\mapsto\Phi_{X_\xi}(t;g)$ exists on all of $\mathbb{R}\times G$.
\end{theorem}
\begin{proof}
We may assume that $\Phi_{X_\xi}(\Delta t;e)$ exists for sufficiently small $\Delta t\in\mathbb{R}$.
Let
\begin{align*}
	\tilde{\Phi}_{X_\xi}(t;g)
	&= (R_g\circ\Phi_{X_\xi})(t;e)
	\ .
\end{align*}
Evidently, $\tilde{\Phi}_{X_\xi}(0;g)=g$ and
\begin{align*}
	\frac{\dt}{\dt t}
		\tilde{\Phi}_{X_\xi}(t;g)
	\Big|_{t=0}
	= \Tt R_g|_e \cdot X_\xi|_e
	= X_\xi|_g
	\ ,
\end{align*}
from which it follows that $\tilde{\Phi}_{X_\xi}(t;g)=\Phi_{X_\xi}(t;g)$.
Thus, the flow starting at $g$ is reduced to the flow starting at the identity element~$e$.
Via the semigroup property, the existence of $\Phi_{X_\xi}(\Delta t;e)$ therefore implies that of $\Phi_{X_\xi}(2\Delta t;e)=\Phi_{X_\xi}(\Delta t;g=\Phi_{X_\xi}(\Delta t;e))$.
It follows that the flow $\Phi_{X_\xi}(t;g)$ of the right-invariant vector field~$X_\xi$ exists for all $g\in G$ and all $t\in\mathbb{R}$.
\end{proof}

\subsubsection{Definition of the Exponential Map}

Theorem~\ref{th:PhiGlobal} guarantees the existence of the flow of a right-invariant vector field for all times.
This guarantees the existence of the map defined below.
\begin{definition}
The (Lie) \mydef{exponential map} is defined as
\begin{align*}
	\exp: 
	\mathfrak{g} &\rightarrow G
	\ ,\nonumber\\
	\xi &\mapsto \exp\xi := \Phi_{X_\xi}(1; e)
	\ .
\end{align*}
Here, $X_\xi$ denotes the right-invariant vector field associated with $\xi$.
\end{definition}

\begin{example}
We consider $\SO(n)$ again.
In Example~\ref{ex:flowSOn}, the flow of a right-invariant vector field $X_A$ associated with the tangent vector $X_A|_I=A\in\fso(n)$ was given by $\Phi_{X_A}:(t,Q)\mapsto\exp_\mathsf{M}(tA)\,Q$.
For $\exp:\fso(n)\rightarrow\SO(n)$, we thus find that
\begin{align*}
	\exp A &= \Phi_{X_A}(1; I) = \exp_\mathsf{M}(A)
	\ .
\end{align*}
Hence, the Lie exponential map on $\SO(n)$ coincides with the matrix exponential $\exp_\mathsf{M}$.
\end{example}

Conversely, using the exponential map, the flow of a \emph{right-invariant} vector field $X_\xi$ with $\xi\in\mathfrak{g}$ can be represented transparently.

\begin{lemma}
Let $\xi\in\fg$ with associated right-invariant vector field $X_\xi$.
Then
\begin{align*}
	\exp(t\xi) &= \Phi_{X_\xi}(t; e)
	\ .
\end{align*}
\end{lemma}
\begin{proof}
For $t\in\mathbb{R}$, let the map $\tilde{\Phi}:\mathbb{R}\rightarrow G$ be given by $\tilde{\Phi}(s)=\Phi_{X_\xi}(st; e)$.
We have $\tilde{\Phi}(0)=e$ and
\begin{align*}
	\frac{\dt}{\dt s}	
	\tilde{\Phi}(s)
	&=
	\frac{\dt}{\dt s}
	\Phi_{X_\xi}(st; e)
	= tX_\xi|_{\Phi_{X_\xi}(st; e)} = tX_\xi|_{\tilde{\Phi}(s)}
	\ ,
\end{align*}
from which it follows that $\tilde{\Phi}(s)=\Phi_{tX_\xi}(s; e)$.
Evaluating at $s=1$ yields the claim.
\end{proof}

\begin{lemma}
The derivative of $t\mapsto\exp(t\xi)$ is given by
\begin{align*}
	\frac{\dt}{\dt t}\exp(t\xi)
	&=
	\Tt R_{\exp(t\xi)}|_e\cdot\xi
	=
	\Psi_1(\exp(t\xi); \xi)
	\ .
\end{align*}
\end{lemma}
\begin{proof}
The statement follows from the preceding lemma via
\begin{align*}
	\frac{\dt}{\dt t}\exp(t\xi)
	&= 
	\frac{\dt}{\dt t}
	\Phi_{X_\xi}(t;e)
	= X_\xi|_{\exp(t\xi)}
	= \Tt R_{\exp(t\xi)}|_e\cdot\xi
	\ .
\end{align*}
\end{proof}

\begin{proposition}
\label{th:PhiFromExp}
For all $(t,g)\in\mathbb{R}\times G$, the flow $\Phi_{X_\xi}$ of a right-invariant vector field $X_\xi$ with $X_\xi|_e=\xi$ can be reduced to the exponential map $\exp$ via
\begin{align*}
	\Phi_{X_\xi}(t;g)
	&= \exp(t\xi)\,g
	\ .
\end{align*}
\end{proposition}
\begin{proof}
Using the preceding lemma, we compute
\begin{align*}
	\frac{\dt}{\dt t}\exp(t\xi)\,g
	&=
	\Tt R_g|_{\exp(t\xi)}\cdot\frac{\dt}{\dt t}\exp(t\xi)
	=
	\Tt R_g|_{\exp(t\xi)}\cdot\Tt R_{\exp(t\xi)}|_e\cdot\xi
	\nonumber\\
	&=
	\Tt R_{\exp(t\xi)g}|_e\cdot\xi = X_\xi|_{\exp(t\xi)g}
	\ ,
\end{align*}
so that $\exp(t\xi)\,g$ satisfies the differential equation of $\Phi_{X_\xi}(t;g)$.

Moreover, $\exp(t\xi)\,g|_{t=0}=\exp(0)\,g$.
The flow of the zero vector field $\Phi_0(t; g)$ is constant for all $g$, since $\dot{\Phi}_0(t; g)=0$. It follows that $\exp(0)=\Phi_0(1; e)=e$, so that the initial value $\exp(t\xi)\,g|_{t=0}=g$ also agrees with that of $\Phi_{X_\xi}(t;g)$.
\end{proof}

\subsubsection{Properties of the Exponential Map}

In this section, we record several analytic and algebraic properties of the exponential map.

\begin{proposition}
For fixed $\xi\in\mathfrak{g}$, the map $t\mapsto\exp(t\xi)$ is a group homomorphism $(\mathbb{R},+)\rightarrow G$.
Its image is a \mydef{one-parameter subgroup} of $G$.
\end{proposition}
\begin{proof}
For $t,s\in\mathbb{R}$, using \ref{th:PhiFromExp}, we have
\begin{align*}
	\exp((t+s)\xi)
	&= \exp((t+s)\xi)e
	= \Phi_{X_\xi}(t+s;e)
	= \Phi_{X_\xi}(t; \Phi_{X_\xi}(s;e))
	\nonumber\\
	&= \exp(t\xi)\Phi_{X_\xi}(s;e) = \exp(t\xi)\exp(s\xi)
	\ ,
\end{align*}
where $X_\xi$ again denotes the right-invariant vector field with $X_\xi|_e=\xi$.
\end{proof}


\begin{lemma}
At the origin $0\in\fg$, the derivative of the exponential map equals the identity,
\begin{align*}
	\Dt\exp|_0
	&=
	\mathrm{id}
	\ .
\end{align*}
\end{lemma}
\begin{proof}
For the exponential map at the origin and for $\xi\in\fg$ with associated right-invariant vector field $X_\xi$, we obtain
\begin{align*}
	\Dt\exp|_0 \cdot\xi
	&=
	\frac{\dt}{\dt t}
		\exp(t\xi)
	\Big|_{t=0}
	= \frac{\dt}{\dt t}
		\Phi_{X_{\xi}}(t;e)
	\Big|_{t=0}
	= X_\xi|_e = \xi
	\ .
\end{align*}
\end{proof}

Since the derivative of $\exp$ at the origin equals the identity,
there exists a neighborhood $U\subset\mathfrak{g}$ on which $\exp$ is a diffeomorphism onto its image $V=\exp U$.

\begin{definition}
In the given situation, we denote the inverse map $V\rightarrow U$ of $\exp$ by
\begin{align*}
	\ln &= \exp^{-1}
	\ .
\end{align*}
The map $\ln$ is a chart map on $V$.
\end{definition}
By composing $\ln\circ R_{g^{-1}}$ with right multiplication, a local chart can be obtained for every $g\in G$.

\begin{example}
Let $a\in\mathbb{R}^3$ be a (coordinate) vector, which we identify with $A\in\fso(3)$ via $A=[a]_\times$.
By $|a|=\sqrt{(a^x)^2+(a^y)^2+(a^z)^2}$ we denote the standard Euclidean norm of $a$.
\mydef{Rodrigues' formula} states that
\begin{align*}
	\exp_\mathsf{M}(A)
	&=
	I + \frac{\sin(|a|)}{|a|}\,A + \frac{1 - \cos(|a|)}{|a|^2}\,A^2
	\ .
\end{align*}
$\exp_\mathsf{M}(A)$ describes a (right-handed) rotation by the angle $|a|$ around the axis $a/|a|$.\footnotemark
Hence, $\exp_\mathsf{M}:\fso(3)\rightarrow \SO(3)$ is in any case surjective.
However, the map is not globally injective because, for instance, for $|a|=\pi$ the values $\exp_\mathsf{M}(A)$ and $\exp_\mathsf{M}(-A)$ coincide.

Since $t\mapsto\exp_\mathsf{M}(tA)$ is an integral curve of the (right-invariant) vector field defined by $A$ and since such integral curves do not intersect, an equality $\exp_\mathsf{M}(A)=\exp_\mathsf{M}(B)$ with $B=[b]_\times\in\fso(3)$ implies that $A$ is a real multiple of $B$.
If $|a|<\pi$ and $|b|<\pi$, the equality of their $\exp_\mathsf{M}$-images even implies $a=b$, as can be seen directly from Rodrigues' formula.

For $a\in\mathrm{B}^3_\pi=\{a\in\mathbb{R}^3\,|\,|a|<\pi\}$, the map $a\mapsto\exp_\mathsf{M}([a]_\times)$ is thus a smooth bijection onto its image.
On $\exp_\mathsf{M}([\mathrm{B}^3_{\pi/2}]_\times)$, we can write down the inverse explicitly as
\begin{align*}
	\ln_\mathsf{M}:
	\exp_\mathsf{M}([\mathrm{B}^3_{\pi/2}]_\times)
	&\rightarrow \fso(3)
	\ ,\nonumber\\
	C 
	&\mapsto
	\frac{\arcsin|k_C|}{|k_C|}
	\, K_C
\end{align*}
with $K_C=[k_C]_\times=(C-C\trans)/2$.
Since $\ln_\mathsf{M}$ is evidently also smooth on its domain, it serves as a chart for all rotations with a rotation angle strictly less than $\pi/2$ in absolute value.

The components of the vectors $a\in\mathbb{R}^3$ considered here are the coordinates with respect to the basis of $\fso(3)$ given in Example~\ref{ex:so3base}.
They are thus normal coordinates of $\SO(3)$ in the sense of Definition~\ref{def:normalcoords}.
\end{example}
\footnotetext{
	For $a=\theta\,e_x$ with $\theta>0$, for instance, we find
	\begin{align*}
		\exp_\mathsf{M}([\theta\,e_x]_\times)
		&=
		\begin{pmatrix}
			1 & 0 & 0 \\
			0 & \cos\theta & -\sin\theta \\
			0 & \sin\theta & \cos\theta
		\end{pmatrix}
		\ .
	\end{align*}		
	For an arbitrary $a$, there exist rotation matrices $Q$ satisfying $Qa=|a|e_x$.
	In this case,
	$\exp_\mathsf{M}([a]_\times)=Q\exp_\mathsf{M}(|a|[e_x]_\times)Q^{-1}$
	is a rotation around the axis $a/|a|$ by the angle $|a|$.
}

\begin{proof}[Proof of Rodrigues' Formula]
We set
\begin{align*}
	F(t,A) &=
	I 
	+ \frac{\sin(t|a|)}{|a|}\,A
	+ \frac{1 - \cos(t|a|)}{|a|^2}\,A^2
\end{align*} 
and show the equality $F(t,A)=\exp_\mathsf{M}(tA)$.

For $A=[a]_\times$, the Grassmann vector triple product identity yields $A^2=aa\trans - |a|^2I$.
From this, it follows in turn that $A^3=-|a|^2A$.
Hence,
\begin{align*}
	A\,F(t,A)
	&=
	\cos(t|a|)\,A
	+
	\frac{\sin(t|a|)}{|a|}\,A^2
	\ .	
\end{align*}
The same expression is obtained for the derivative $\dt F(t,A)/\dt t$.

Thus, $F(t,A)$ satisfies the same differential equation as $\exp_\mathsf{M}(tA)$ and, moreover, has the same initial value $F(0,A)=I$.
\end{proof}

\begin{definition}
\label{def:normalcoords}
Let $\xi_1,\ldots,\xi_n$ be a basis of the Lie algebra $\mathfrak{g}$ of the Lie group $G$.
The coordinates $x^i\in\mathbb{R}$, $i=1,\ldots,n$, used in the coordinate representation $\xi=\sum_{i=1}^n x^i\xi_i$ for $\xi\in\mathfrak{g}$ are called \mydef{(Lie) normal coordinates}.
\end{definition}
Via $\exp$, the Lie normal coordinates provide a (local) parametrization of $G$.

\subsubsection{Exponentialabbildung und adjungierte Darstellung}

Das Diagramm
\begin{center}
\begin{tikzcd}
\mathfrak{g} \arrow{d}{\exp} \arrow{r}{\ad}
&\mathfrak{gl}(\mathfrak{g}) \arrow{d}{\exp}\\
G \arrow{r}{\Ad}
& \mathrm{GL}(\mathfrak{g})
\end{tikzcd}
\end{center}
illustriert zwei Wege für Abbildungen $\mathfrak{g}\rightarrow\mathrm{GL}(\mathrm{g})$.
Dabei ist $\exp:\mathfrak{gl}(\mathfrak{g})\rightarrow\mathrm{GL}(\mathfrak{g})$ die ``gewöhnliche'' Operatorexponentialfunktion.

Dass sie übereinstimmen, dass also
\begin{empheq}[box=\resultbox]{align}
	\Ad\circ\exp &= \exp\circ\ad
\end{empheq}
ist, soll hier gezeigt werden.

Sei dazu $\xi\in\mathfrak{g}$ beliebig.
Wir betrachten zwei Kurven in $\mathrm{GL}(\mathfrak{g})$,
\begin{align*}
	\Gamma_1: t&\mapsto \exp(t\,\ad(\xi))
	\ ,\\
	\Gamma_2: t&\mapsto \Ad(\exp(t\xi))
	\ .
\end{align*}
$\Gamma_1$ ist nach Definition eine Einparametergruppe.
Dass dies auch $\Gamma_2$ ist, folgt aus
\begin{align*}
	\Gamma_2(t+s)
	&= \Ad(\exp((t+s)\xi))
	= \Ad(\exp(t\xi)\exp(s\xi))
	\nonumber\\
	&= \Ad(\exp(t\xi)\cdot\Ad(\exp(s\xi)
	\ .
\end{align*}
Es ist außerdem
\begin{align*}
	\dot{\Gamma}_1(0)
	&= \ad(\xi)
\end{align*}
und
\begin{align*}
	\dot{\Gamma}_2(0)
	&= \Dt\,\Ad|_e\cdot \xi = \ad(\xi) = \dot{\Gamma}_1(0)
	\ .
\end{align*}
Aufgrund der Einparametergruppeneigenschaft gilt für $i=1,2$
\begin{align*}
	\dot{\Gamma}_i(t)
	&=
	\frac{\dt}{\dt s}
	\Gamma_i(s+t)
	\Big|_{s=0}
	= \dot{\Gamma}_i(0)\Gamma_i(t)
	= \ad(\xi)\circ\Gamma_i(t)
\end{align*}
(wobei $\Gamma_i(t)$ und $\ad(\xi)$ vermöge der Einparametergruppe kommutieren).

Beide Kurven erfüllen also dieselbe gewöhnliche Differenzialgleichung ersten Grades und haben wegen $\Gamma_1(0)=\Gamma_2(0)=\mathrm{id}$ überdies den gleichen Anfangswert.
Damit stimmen $\Gamma_1$ und $\Gamma_2$ überein, 
so dass die Gleichheit $\Gamma_1(1)=\Gamma_2(1)$ die zu zeigende Aussage ergibt.

\subsubsection{Ableitung der Exponentialfunktion}

Ausgehend von Gl.~(\ref{eq:dExp0}) betrachten wir die Ableitung mit um den Faktor $s$ skalierten Tangentialvektoren,
\begin{align*}
	\Dt\exp|_{s\xi} \cdot s\upsilon
	&=
	\Tt R_{\exp(s\xi)}^{-1}|_{\exp(s\xi)}
	\cdot
	\frac{\dt}{\dt t}
		\exp(s\xi + st\upsilon)
	\Big|_{t=0}
	\nonumber\\
	&=
	\frac{\dt}{\dt t}
		\exp(s\xi + st\upsilon)
		\exp(-s\xi)
	\Big|_{t=0}
\end{align*}
Mit Hilfe der durch 
\begin{align*}
	h_{\xi,\upsilon}(t, s)
	&:= 
	\exp(s\xi + st\upsilon)\exp(-s\xi)
\end{align*}
definierten Abbildung $h_{\xi,\upsilon}:\mathbb{R}^2\rightarrow G$ schreiben wir
\begin{align*}
	\Dt\exp|_{s\xi} \cdot s\upsilon
	&= \dot{h}_{\xi, \upsilon}(0, s)
	\ ,
\end{align*}
wobei
\begin{align*}
	\dot{h}_{\xi, \upsilon}(t, s)
	&= \frac{\dt}{\dt t} h_{\xi, \upsilon}(t, s)
	\ .
\end{align*}
Wir notieren einige in der Folge wichtige Eigenschaften dieser Abbildung,
\begin{align*}
	h_{\xi,\upsilon}(0, s)
	&= e
	\ ,\\
	\dot{h}_{\xi, \upsilon}(0, s) 
	&\in\mathfrak{g}
	\ ,\\
	\dot{h}_{\xi, \upsilon}(0, 0)
	&= 0
	\ ,\\
	\dot{h}_{\xi, \upsilon}(0, 1)
	&= \Dt\exp|_{\xi} \cdot \upsilon
	\ .
\end{align*}
Die explizite Berechnung von $\dot{h}_{\xi, \upsilon}(0, s)$ erfolgt über die Ableitung nach $s$, beginnend mit der Produktregel (für $\Tt\,\mathrm{mult}$),
\begin{align*}
	\frac{\partial}{\partial s}
	\dot{h}_{\xi, \upsilon}(0, s)	
	&=
	\frac{\dt}{\dt t}	
	\Big(
		\Tt R_{\exp(-s\xi)}
		|_{\exp(s\xi + st\upsilon)}
		\cdot	
		\frac{\partial}{\partial s}\exp(s\xi + st\upsilon)
	\nonumber\\
	&\qquad\qquad+
		\Tt L_{\exp(s\xi + st\upsilon)}
		|_{\exp(-s\xi)}		
		\cdot			
		\frac{\partial}{\partial s}
		\exp(-s\xi)
	\Big)	
	\Big|_{t=0}
	\nonumber\\	
	&=
	\frac{\dt}{\dt t}	
	\Big(
		\Tt R_{\exp(-s\xi)}
		|_{\exp(s\xi + st\upsilon)}
		\cdot	
		\Tt R_{\exp(s\xi + st\upsilon)}
		|_e
		\cdot
		(\xi + t\upsilon)
	\nonumber\\
	&\qquad\qquad-
		\Tt L_{\exp(s\xi + st\upsilon)}
		|_{\exp(-s\xi)}	
		\cdot
		\Tt R_{\exp(-s\xi)}
		|_e
		\cdot
		\xi
	\Big)	
	\Big|_{t=0}
	\nonumber\\	
	&=
	\frac{\dt}{\dt t}	
	\Big(
		\Tt R_{h_{\xi, \upsilon}(t, s)}
		|_e
		\cdot
		(\xi + t\upsilon)
	\nonumber\\
	&\qquad\qquad-
		\Tt L_{h_{\xi, \upsilon}(t, s)}
		|_{e}	
		\cdot
		\Ad(\exp(-s\xi))
		\cdot
		\xi
	\Big)	
	\Big|_{t=0}
	\ .
\end{align*}
Mit
\begin{align*}
	\frac{\dt}{\dt t}	
	\Big(
		\Tt R_{h_{\xi, \upsilon}(t, s)}
		|_e
		\cdot
		t\upsilon
	\Big)
	\Big|_{t=0}
	&=
	\Tt R_{h_{\xi, \upsilon}(0, s)}
	|_e
	\cdot
	\upsilon
	= \upsilon
\end{align*}
und
\begin{align*}
	\Ad(\exp(-s\xi))\cdot\xi
	&= \exp(\ad(-s\xi))\cdot\xi
	= \xi
\end{align*}
erhalten wir 
\begin{align*}
	\frac{\partial}{\partial s}
	\dot{h}_{\xi, \upsilon}(0, s)	
	&=
	\upsilon +
	\frac{\dt}{\dt t}		
	\Big(
		\Tt R_{h_{\xi, \upsilon}(t, s)}
		|_e -
		\Tt L_{h_{\xi, \upsilon}(t, s)}
		|_{e}
	\Big)
	\Big|_{t=0}
	\cdot\xi
	\ .
\end{align*}
Mit
\begin{align*}
	\big(
	\Tt R_{h_{\xi, \upsilon}(t, s)}
	|_e -
	\Tt L_{h_{\xi, \upsilon}(t, s)}
	|_{e}
	\big)\big|_{t=0}
	&= 0
\end{align*}
sowie $\Tt L_{h_{\xi, \upsilon}(0, s)^{-1}}|_{h_{\xi, \upsilon}(0, s)}=\mathrm{id}$ (was trivial ist) folgt
\begin{align*}
	&\frac{\dt}{\dt t}	
	\Big(
		\Tt R_{h_{\xi, \upsilon}(t, s)}
		|_e -
		\Tt L_{h_{\xi, \upsilon}(t, s)}
		|_{e}	
	\Big)	
	\Big|_{t=0}
	\nonumber\\
	&=
	\frac{\dt}{\dt t}
	\big(
		\Tt L_{h_{\xi, \upsilon}(t, s)^{-1}}
		|_{h_{\xi, \upsilon}(t, s)} \cdot
		\Tt R_{h_{\xi, \upsilon}(t, s)}
		|_e -
		\mathrm{id}	
	\big)
	\Big|_{t=0}
	\nonumber\\
	&=
	\frac{\dt}{\dt t}
		\Ad(h_{\xi, \upsilon}(t, s)^{-1})
	\Big|_{t=0}
	= \ad(-\dot{h}_{\xi, \upsilon}(0, s))
	\ .
\end{align*}
Zusammengefasst erhalten wir
\begin{align*}
	\frac{\partial}{\partial s}
	\dot{h}_{\xi, \upsilon}(0, s)
	&=
	\ad(-\dot{h}_{\xi, \upsilon}(0, s))\cdot\xi + \upsilon
	= \ad(\xi)\cdot\dot{h}_{\xi, \upsilon}(0, s) + \upsilon
	\ .
\end{align*}
Dies ist eine gewöhnliche Diffenzialgleichung für $\dot{h}_{\xi, \upsilon}(0, s)$ auf dem endlichdimensionalen Vektorraum $\mathfrak{g}$.
Mit dem Anfangswert $\dot{h}_{\xi, \upsilon}(0, 0)=0$ finden wir
\begin{align*}
	\dot{h}_{\xi, \upsilon}(0, s)
	&=
	\int_0^s
	\exp((s-s')\,\ad(\xi))\cdot{\upsilon}
	\ \dt s'
\end{align*}
und mit $\dot{h}_{\xi, \upsilon}(0, 1)=\Dt\exp|_{\xi}\cdot\upsilon$ schließlich
\begin{empheq}[box=\resultbox]{align}
	\Dt\exp|_{\xi}
	&= E(\ad(\xi))
	\ ,
\end{empheq}
wobei
\begin{align}
	E(x) &:=
	\int_0^1
	\mathrm{e}^{(1-s')x}\ \dt s'
	= \frac{\mathrm{e}^x-1}{x}
	= 1 + \frac{1}{2}x + \frac{1}{6}x^2 + \ldots
	\ .
\end{align}

\newpage

\section{Zusammenhänge auf Lie-Gruppen}

Dieser Teil ist der Diskussion von Zusammenhängen auf Lie-Gruppen gewidmet,
die aus Sicht mancher Anwendung, insbesondere der Navigation, interessant sind.
Ein allgemeiner Überblick über Zusammenhänge auf Lie-Gruppen ist nicht Ziel.
Auch beschränken wir uns im Detail auf affine Zusammenhänge auf dem (parallelisierten) Tangentialbündel.

\subsection{Parallelisierung von $\Tt G$ und $\Tt\Tt G$}

\subsubsection{$\Tt G$ als Lie-Gruppe}

Das Tangentialbündel $\Tt G$ einer Lie-Gruppe~$G$ ist vermöge der Multiplikation
\begin{align}
	\mathrm{mult}_{\Tt G}:
	\Tt G\times\Tt G 
	&\rightarrow \Tt G
	\ ,\nonumber\\
	(\xi_g, \upsilon_h)
	&\mapsto
	\Tt\,\mathrm{mult}|_{(g,h)}
	\cdot
	(\xi_g, \upsilon_h)
\end{align}
und Inversenbildung
\begin{align}
	\mathrm{inv}_{\Tt G}:
	\Tt G
	&\rightarrow \Tt G
	\ ,\nonumber\\
	\xi_g
	&\mapsto
	\Tt\,\mathrm{inv}|_{g}
	\cdot
	\xi_g
\end{align}
selbst eine Lie-Gruppe.
Dabei ist der Nullvektor $0_e=0\in\Tt_eG=\mathfrak{g}$ am neutralen Element von $G$ das neutrale Element von $\Tt G$.

Die Abbildung $\Psi_1^{-1}:\xi_g\mapsto\Tt R_{g^{-1}}|_g\cdot \xi_g$, vgl. Defitinion~\ref{def:Psi1}, definiert den Diffeomorphimus $\Tt G\rightarrow G\times\mathfrak{g}$, den wir zur Parallelisierung von $G$ verwenden.
Über diesen transportieren wir die Lie-Gruppenstruktur auf $G\times\mathfrak{g}$ und finden mit
\begin{align*}
	&\Tt R_{(gh)^{-1}}|_{gh}
	\cdot
	\Tt\mathrm{mult}|_{(g,h)}
	\cdot
	\begin{pmatrix}
		\Tt R_g|_e\cdot \xi \\
		\Tt R_h|_e\cdot \upsilon
	\end{pmatrix}
	\nonumber\\
	&=
	\Tt(R_{(gh)^{-1}}\circ R_h\circ R_g)|_e\cdot\xi
	+
	\Tt(R_{(gh)^{-1}}\circ L_g\circ R_h)|_e\cdot\xi
	\nonumber\\
	&= \xi + \Ad(g)\cdot\upsilon 
\end{align*}
für die Multiplikation
\begin{align}
	\mathrm{mult}_{G\ltimes_{\Ad}\mathfrak{g}}:
	(G\times\mathfrak{g})\times(G\times\mathfrak{g})
	&\rightarrow G\times\mathfrak{g}
	\nonumber\\
	((g,\xi), (h,\upsilon))
	&\mapsto
	(gh, \xi + \Ad(g)\cdot\upsilon)
	\ .
\end{align}
Die Lie-Gruppenstruktur ist damit jene des semidirekten Produkts
$G\ltimes_{\Ad}\mathfrak{g}$ mittels des Homomorphismus $\Ad:G\rightarrow\mathrm{GL}(\mathfrak{g})$.
Ihr neutrales Element ist $(e,0)$, 
das inverse Element zu $(g,\xi)$ ist offenbar $(g^{-1},-\Ad(g^{-1})\cdot\xi)$.

Die Abbildung~$\Psi_1:G\ltimes_{\Ad}\mathfrak{g}\rightarrow\Tt G$ ist somit nach Konstruktion ein Lie-Gruppen\-isomorphismus.

\subsubsection{Lie-Algebra von $\Tt G$}

Vermöge $\Psi_1$ können wir die Lie-Algebrastruktur von $\Tt G$ statt auf $\Tt G$ im semidirekten Produkt $G\ltimes_{\Ad}\mathfrak{g}$ studieren.
Dazu führen wir zunächst die Rechtsmultiplikation\footnote{
	Für die Rechtsmultiplikation auf $G\ltimes_{\Ad}\mathfrak{g}$ und später auch $\Tt G$ verwenden wir ebenfalls das Symbol ``$R$''.
	Durch das Argument im Index, hier $(g,\xi)$, sollte aber jeweils klar sein, in welchem Raum $R$ operiert.
}
ein.
Sei dazu $(g,\xi)\in G\ltimes_{\Ad}\mathfrak{g}$ und
\begin{align}
	R_{(g,\xi)}=\mathrm{mult}_{G\ltimes_{\Ad}\mathfrak{g}}
	(-, (g,\xi)):
	G\ltimes_{\Ad}\mathfrak{g} &\rightarrow G\ltimes_{\Ad}\mathfrak{g}
	\ .
\end{align}
Sei $\Xi|_{(e,0)}=(\Xi_1, \Xi_2)\in\mathfrak{g}\times\mathfrak{g}$ und sei $\Gamma=(\Gamma_1,\Gamma_2):(-\varepsilon,\varepsilon)\rightarrow G\ltimes_{\Ad}\mathfrak{g}$ eine Kurve mit $\Gamma_1(0)=h$, $\Gamma_2(0)=\upsilon$, $\dot{\Gamma}_1(0)=\Tt R_h|_e\cdot\Xi_1$ und $\dot{\Gamma}_2(0)=\Xi_2$.
Dann ist
\begin{align*}
	\frac{\dt}{\dt t}
	\big(
		R_{(g,\xi)}\cdot\Gamma(t)
	\big)
	\Big|_{t=0}
	&=
	\frac{\dt}{\dt t}
	\begin{pmatrix}
		\Gamma_1(t)\,g \\
		\Gamma_2(t) + \Ad(\Gamma_1(t))\cdot\xi
	\end{pmatrix}
	\Big|_{t=0}
	\nonumber\\
	&=
	\begin{pmatrix}
		\Tt R_g|_h\cdot\Tt R_h|_e\cdot\Xi_1 \\
		\Xi_2 + \Dt_{\Xi_1}\Ad|_{h}\cdot\xi
	\end{pmatrix}
	\ .
\end{align*}
Durch Auswerten bei $(h,\upsilon)=(e,0)$ finden wir für das durch $\Xi|_{(e,0)}$ definierte rechtsinvariante Feld $\Xi$ somit
\begin{empheq}[box=\resultbox]{align}
	\Xi|_{(g,\xi)}
	&=
	\Tt R_{(g,\xi)}|_{(e,0)}
	\cdot
	\Xi|_{(e,0)}
	=
	\begin{pmatrix}
		\Tt R_g|_e\cdot\Xi_1 \\
		\Xi_2 + \ad(\Xi_1)\cdot\xi
	\end{pmatrix}
	\in
	\Tt_gG\times\mathfrak{g}
	\ ,
\end{empheq}
wobei $\Tt_\xi\mathfrak{g}$ und $\mathfrak{g}$ miteinander identifiziert werden.

Für die Linksmultiplikation finden wir analog das Differenzial
\begin{align*}
	\Tt L_{(g,\xi)}|_{(h,\upsilon)}
	\cdot
	\dot{\Gamma}(0)
	&=
	\frac{\dt}{\dt t}
	\big(
		L_{(g,\xi)}\cdot\Gamma(t)
	\big)
	\Big|_{t=0}
	=
	\frac{\dt}{\dt t}
	\begin{pmatrix}
		g\,\Gamma_1(t) \\
		\xi + \Ad(g)\cdot\Gamma_2(t)
	\end{pmatrix}
	\Big|_{t=0}
	\nonumber\\
	&=	
	\begin{pmatrix}
		\Tt L_g|_h\cdot\dot{\Gamma}_1(0) \\
		\Ad(g)\cdot{\Gamma}_2(0)
	\end{pmatrix}
	\ ,
\end{align*}
also $\Tt L_{(g,\xi)}|_{(h,\upsilon)}=\Tt L_g|_h\oplus\Ad(g)$.

Sei im Folgenden $\Gamma$ eine \emph{Integralkurve} zum rechtsinvarianten Vektorfeld $\Xi$, also $\dot{\Gamma}(t)=\Xi|_{\Gamma(t)}$, 
mit $\Gamma(0)=(e,0)$
Sei ferner $\Upsilon$ ein weiteres durch $\Upsilon|_{(e,0)}=(\Upsilon_1, \Upsilon_2)\in\mathfrak{g}\times\mathfrak{g}$ definiertes rechtsinvariantes Vektorfeld.
Damit berechnen wir die Lie-Klammer gemäß
\begin{align*}
	\big[
		\Xi|_{(e,0},\Upsilon|_{(e,0}
	\big]^{\mathfrak{g}\times\mathfrak{g}}
	&= 
	\frac{\dt}{\dt t}
	\big(
		[L_{\Gamma(t)}]_*
		\Psi
	\big)|_{(e,0)}
	\Big|_{t=0}
	\ .
\end{align*}
Mit
\begin{align*}
	\big(
		[L_{\Gamma(t)}]_*
		\Psi
	\big)|_{(e,0)}
	&=
	\Tt	L_{\Gamma(t)}|_{\Gamma(t)^{-1}}
	\cdot
	\Tt R_{\Gamma(t)^{-1}}|_{(e,0)}
	\cdot
	\begin{pmatrix}
		\Upsilon_1 \\ \Upsilon_2
	\end{pmatrix}
	\nonumber\\
	&=
	\begin{pmatrix}
		\Ad(\Gamma_1(t))
		\cdot\Upsilon_1
		\\
		\Ad(\Gamma_1(t))\cdot\Upsilon_2
		-
		\Ad(\Gamma_1(t))
		\cdot\ad(\Upsilon_1)
		\cdot\Ad(\Gamma_1(t))^{-1}
		\cdot\Gamma_2(t)
	\end{pmatrix}
\end{align*}
finden wir schließlich
\begin{empheq}[box=\resultbox]{align}
	\Big[
		\begin{pmatrix}
			\Xi_1 \\ \Xi_2
		\end{pmatrix}
		,
		\begin{pmatrix}
			\Upsilon_1 \\ \Upsilon_2
		\end{pmatrix}
	\Big]^{\mathfrak{g}\times\mathfrak{g}}
	&=
	\begin{pmatrix}
		 \ad(\Xi_1)\cdot\Upsilon_1
		 \\
		 \ad(\Xi_1)\cdot\Upsilon_2
		 -
 		 \ad(\Upsilon_1)\cdot\Xi_2
	\end{pmatrix}	
	\nonumber\\
	&=
	\begin{pmatrix}
		[\Xi_1,\Upsilon_1]^\mathfrak{g}
		 \\
		 [\Xi_1,\Upsilon_2]^\mathfrak{g}
		 -
		 [\Upsilon_1,\Xi_2]^\mathfrak{g}
	\end{pmatrix}	
	\ .
\end{empheq}

\subsubsection{Parallelisierung von $\Tt\Tt G$}

Sei $(h,\upsilon)\in G\times\mathfrak{g}$ und sei $(\Xi_1,\Xi_2)\in \mathfrak{g}\times\mathfrak{g}$.
Wir betrachten eine glatte Kurve $\Gamma=(\Gamma_1,\Gamma_2):(-\varepsilon,\varepsilon)\rightarrow G\times\mathfrak{g}$ mit Startpunkt $\Gamma(0)=(h,\upsilon)$ und Startgeschwindigkeit $\dot{\Gamma}(0)=(\Tt R_h|_e\cdot\Xi_1,\Xi_2)$.\footnote{
	Wir identifizieren den Vektorraum $\Tt_\upsilon\mathfrak{g}$ mit dem Vektorraum $\mathfrak{g}$ selbst.
	--- Für die Kurve können wir beispielsweise $\Gamma(t)=(\exp(t\Xi_1)h,\upsilon + t\Xi_2)$ wählen.
}

Ableiten der Parallelisierung $\Psi_1:G\times\mathfrak{g}\rightarrow\Tt G$ längs der betrachteten Kurve ergibt
\begin{align*}
	\frac{\dt}{\dt t}
		\Psi_1(\Gamma(t))
	\Big|_{t=0}
	&=
	\Tt\Psi_1
	\big|_{(h,\upsilon)}
	\cdot
	\begin{pmatrix}
		\Tt R_h|_e\cdot\Xi_1 \\
		\Xi_2
	\end{pmatrix}
	\ ,
\end{align*}
was die Abbildung
\begin{empheq}[box=\definitionbox]{align}
	\Psi_2:
	(G\times\mathfrak{g})\times(\mathfrak{g}\times\mathfrak{g})
	&\rightarrow \Tt\Tt G
	\ ,\nonumber\\
	(h,\upsilon\,;\,\Xi_1,\Xi_2)
	&\mapsto
	\Tt\Psi_1
	\big|_{(h,\upsilon)}
	\cdot
	\begin{pmatrix}
		\Tt R_h|_e\cdot\Xi_1 \\
		\Xi_2
	\end{pmatrix}	
\end{empheq}
motiviert.
$\Psi_2$ ist konstruktionsbedingt offenbar ein Diffeomorphismus.

Sei
\begin{empheq}[box=\definitionbox]{align}
	\psi_2=\Psi_2(e,0;-,-): \mathfrak{g}\times\mathfrak{g}
	&\rightarrow \Tt_{0_e}\Tt G
	\ .
\end{empheq}
Die Abbildung $\psi_2$ ist ein Lie-Algebraisomorphismus, wenn $\mathfrak{g}\times\mathfrak{g}$ die Lie-Klammer $[-,-]^{\mathfrak{g}\times\mathfrak{g}}$ trägt. 
Wegen $\Psi_2(e, 0;\Xi_1,\Xi_2)=\Tt\Psi_1|_{(e,0)}\cdot(\Xi_1,\Xi_2)$ ist nämlich
\begin{empheq}[box=\resultbox]{align}
	\psi_2
	&= \Tt\Psi_1
	\big|_{(e,0)}
\end{empheq}
und $\Psi_1$ ein Lie-Gruppenisomorphimus.

Die Abbildung~$\Psi_2$ lässt sich aus $\psi_2$ mit Hilfe der Rechtsmultiplikation auf $\Tt G$,
\begin{align}
	R_{\upsilon_h}:
	\Tt G &\rightarrow \Tt G
	\ ,\nonumber\\
	\xi_g
	&\mapsto 
	\mathrm{mult}_{\Tt G}(\xi_g, \upsilon_h)
	\ ,
\end{align}
wobei $\upsilon_h\in\Tt_hG$ und $\xi_g\in\Tt_gG$, rekonstruieren.

Wir beobachten zunächst, dass
\begin{align*}
	R_{\upsilon_h}\Psi_1(g, \xi)
	&= \xi_g\upsilon_h
	= \Psi_1(gh, \xi + \Ad(g)\cdot\upsilon)
\end{align*}
ist.
Sei $\Gamma=(\Gamma_1,\Gamma_2):(-\varepsilon,\varepsilon)\rightarrow G\ltimes_{\Ad}\mathfrak{g}$ eine Kurve mit $\Gamma(0)=(e,0)$ und $\dot{\Gamma}(0)=(\Xi_1,\Xi_2)$.
Wir finden einerseits
\begin{align*}
	\frac{\dt}{\dt t}
		R_{\upsilon_h}\Psi_1(\Gamma(t))
	\Big|_{t=0}
	&=
	\Tt R_{\upsilon_h}|_{(e,0)}
	\cdot
	\Psi_2(e,0\,;\, \Xi_1, \Xi_2)
	\ ,
\end{align*}
andererseits
\begin{align*}
	\frac{\dt}{\dt t}
		\Psi_1(
			\Gamma_1(t)\,h, 
			\Gamma_2(t) + \Ad(\Gamma_1(t))\cdot\upsilon
		)
	\Big|_{t=0}
	&=
	\Tt\Psi_1
	\big|_{(h,\upsilon)}
	\cdot
	\begin{pmatrix}
		\Tt R_h|_e\cdot\Xi_1 \\
		\Xi_2
		+ \ad(\Xi_1)\cdot\upsilon
	\end{pmatrix}
	\nonumber\\
	&=
	\Psi_2(
		h, \upsilon
		\,;\,
		\Xi_1, \Xi_2 + \ad(\Xi_1)\cdot\upsilon
	)
	\ .
\end{align*}
Zusammengefasst finden wir
\begin{empheq}[box=\resultbox]{align}
	\Psi_2(
		h, \upsilon
		\,;\,
		\Xi_1, \Xi_2
	)
	&=
	\Tt R_{\upsilon_h}|_{(e,0)}
	\cdot
	\psi_2(
		\Xi_1, \Xi_2 - \ad(\Xi_1)\cdot\upsilon
	)
	\ .
\end{empheq}


\subsection{Zusammenhänge auf $G$}

\subsubsection{Vertikalraum}

In der Parallelisierung $(G\times\mathfrak{g})\times(\mathfrak{g}\times\mathfrak{g})$ können der zweite und vierte Faktor, wie bereits zuvor getan, identifiziert werden.
In der Tat ist der vierte Faktor der trivial parallelisierte Tangentialraum an einem Punkt~$\upsilon$ des zweiten Faktors, eines Vektorraums.\footnote{
	Formal kann etwa die Kurve 
\begin{align*}
	H=(H_1,H_2): \mathbb{R} &\rightarrow G\times\mathfrak{g}
	\ ,\nonumber\\
	t&\mapsto(g, t\zeta)
\end{align*}
betrachtet werden.
Sie identifiziert $H(1)=(g,\zeta)\in G\times\mathfrak{g}$ und $\dot{H}(0)=(0,\zeta)\in\mathfrak{g}\times\mathfrak{g}$ miteinander.
In praktischen Rechnungen werden Vektoren (bei gleicher Basis) einfach miteinander verrechnet.
}

Sei $\Gamma=(\Gamma_1,\Gamma_2):(-\varepsilon,\varepsilon)\rightarrow G\times\mathfrak{g}$ wieder eine Kurve mit mit Startpunkt $\Gamma(0)=(h,\upsilon)$ und Startgeschwindigkeit $\dot{\Gamma}(0)=(\Tt R_h|_e\cdot\Xi_1,\Xi_2)$.
Mit $\pi$ bezeichnen wir die Bündelprojektion $\Tt G\rightarrow G$ und betrachten die Komposition
\begin{center}
\begin{tikzcd}
\Pi:\ 
(-\varepsilon,\varepsilon)
\arrow{r}{\Gamma}
& G\times\mathfrak{g}
\arrow{r}{\Psi_1}
& \Tt G
\arrow{r}{\pi}
& G
\ .
\end{tikzcd}
\end{center}
Offenbar ist $\Pi(t)=\Gamma_1(t)$ und damit
\begin{align*}
	\dot{\Pi}(0)
	&= \Psi_1(h,\Xi_1)
	\ .
\end{align*}
Andererseits ist nach der Kettenregel
\begin{align*}
	\dot{\Pi}(0)
	&= 
	\Tt\pi
	\big|_{\Psi_1(h,\upsilon)}
	\cdot
	\Tt\Psi_1
	\big|_{(h,\upsilon)}
	\cdot
	\dot{\Gamma}(0)
	= \Tt\pi
	\big|_{\Psi_1(h,\upsilon)}
	\cdot
	\Psi_2(h,\upsilon\,;\,\Xi_1,\Xi_2)
	\ .
\end{align*}
Wir finden somit für die Ableitung der Bündelprojektion die Formel
\begin{empheq}[box=\resultbox]{align}
	\Psi_1^{-1}\circ\Tt\pi\circ\Psi_2:
	(h,\upsilon\,;\,\Xi_1,\Xi_2) \mapsto (h,\Xi_1)
	\ .
\end{empheq}
Für festes $(h,\upsilon)$ ist damit der Kern von $\Psi_1^{-1}\circ\Tt\pi\circ\Psi_2$ der zweite Faktor in $\mathfrak{g}\times\mathfrak{g}=\psi_2^{-1}(\Tt_{e,0}\Tt G)$.
Somit beschreibt das (triviale) Unterbündel $(G\times\mathfrak{g})\times(\{0\}\times\mathfrak{g})$ von $(G\times\mathfrak{g})\times(\mathfrak{g}\times\mathfrak{g})$ in der durch $\Psi_1$ und $\Psi_2$ gegebenen Parallelisierung den \mydef{Vertikalraum} über $\Tt G$ bezüglich der Submersion~$\pi$.

Vektoren im Vertikalraum lassen sich, wie eingangs erwähnt, unmittelbar im Vektorraum $\mathfrak{g}=\Tt_eG$ verrechnen.
Um einem Tangentialvektor $\psi_2(\Xi_1,\Xi_2)\in\Tt_{(e,0)}\Tt G$ im Tangentialraum $\Tt_eG$ ``wiederzuverwerten'', ist also eine Projektion auf den Vertikalraum festzulegen.

\subsubsection{Ableitung von Vektorfeldern}

Ein Vektorfeld $Y$ auf einer Lie-Gruppe $G$ ist ein Schnitt im Tangentialbündel, d.~h. eine Abbildung $Y:G\rightarrow\Tt G$ mit der Eigenschaft $\pi\circ Y=\mathrm{id}$.
Als (glatte) Abbildung zwischen den Mannigfaltigkeiten $G$ und $\Tt G$ besitzt ein Vektorfeld stets die Ableitung
\begin{align*}
	\Tt Y: \Tt G\rightarrow \Tt\Tt G
	\ .
\end{align*}
Nützlicher als diese Ableitung sind in der Regel durch Zusammenhänge definierte kovariante Ableitungen, deren Ergebnisse wieder Vektorfelder auf $G$ sind.
Statt Zusammenhänge direkt durch Rechenvorschriften einzuführen, betrachten wir zunächst die Ableitung $\Tt Y$.

Ein Vektorfeld $Y$ können wir über die Parallelisierung mit einer Abbildung $G\rightarrow\mathfrak{g}$ identifizieren.
Wir definieren das $Y$ eindeutig zugeordnete \mydef{trivialisierte Feld}
\begin{empheq}[box=\definitionbox]{align}
	Y^\mathfrak{g}: 
	G &\rightarrow \mathfrak{g}
	\ ,\nonumber\\
	g &\mapsto 
	\Tt R_{g^{-1}}|_g \cdot Y|_g
	\ ,
\end{empheq}
wobei der Superskript $\mathfrak{g}$ andeutet, dass $Y^\mathfrak{g}|_h\in\mathfrak{g}$ ist.
Neben $Y^\mathfrak{g}$ verwenden wir auch den Graphen
\begin{empheq}[box=\definitionbox]{align}
	\tilde{Y}^\mathfrak{g}: 
	G &\rightarrow G\times\mathfrak{g}
	\ ,\nonumber\\
	g &\mapsto 
	(g, Y^\mathfrak{g}|_g)
	\ .
\end{empheq}

Sei $\gamma:(-\varepsilon,\varepsilon)\rightarrow$ eine glatte Kurve in $G$ mit $\gamma(0)=g$ und $\dot{\gamma}(0)=\Psi_1(g,\xi)=\Tt R_g|_e\cdot\xi$ mit $g\in G$ und $\xi\in\mathfrak{g}$.
Wir betrachten die Komposition
\begin{center}
\begin{tikzcd}
Y\circ\gamma:\ 
(-\varepsilon,\varepsilon)
\arrow{r}{\gamma}
& G
\arrow{r}{\tilde{Y}^\mathfrak{g}}
& G\times\mathfrak{g}
\arrow{r}{\Psi_1}
& \Tt G
\end{tikzcd}
\end{center}
und berechnen
\begin{align*}
	\frac{\dt}{\dt t}
		(Y\circ\gamma)(t)
	\Big|_{t=0}
	&=
	\Tt\Psi_1
	|_{(g, Y^\mathfrak{g}|_{g})}
	\cdot
	\frac{\dt}{\dt t}
		(\tilde{Y}^\mathfrak{g}
		\circ\gamma)(t)
	\Big|_{t=0}
	\nonumber\\
	&=
	\Tt\Psi_1
	|_{(g, Y^\mathfrak{g}|_{g})}
	\cdot
	\begin{pmatrix}
		\Tt R_g|_e\cdot\xi\\
		\Dt Y^\mathfrak{g}|_{g}
		\cdot
		\Tt R_g|_e\cdot\xi
	\end{pmatrix}
	\ .
\end{align*}
Durch Vergleich mit
\begin{align*}
	\Psi_2(h,\upsilon\,;\,\Xi_1,\Xi_2)
	&=
	\Tt\Psi_1
	\big|_{(h,\upsilon)}
	\cdot
	\begin{pmatrix}
		\Tt R_{h}|_e\cdot\Xi_1 \\
		\Xi_2
	\end{pmatrix}		
\end{align*}
finden wir in den Trivialisierungen
\begin{empheq}[box=\resultbox]{align}
\label{eq:TX}
	\Psi_2^{-1}\circ\Tt Y|_g\circ\Psi_1(g,-):
	\mathfrak{g}
	&\rightarrow
	(\{g\}\times\{Y^\mathfrak{g}|_{g}\})
	\times(\mathfrak{g}\times\mathfrak{g})
	\ ,\nonumber\\
	\upsilon
	&\mapsto
	(
	g, Y^\mathfrak{g}|_{g}
	\,;\,
	\xi,
	\Dt Y^\mathfrak{g}|_{g}
	\cdot
	\Tt R_g|_e\cdot\xi
	)
	\ .
\end{empheq}

\subsubsection{Affine Zusammenhänge}

Um aus Formel (\ref{eq:TX}) eine kovariante Ableitung zu gewinnen, bedarf es lediglich einer Projektion auf den Vertikalraum.
Die allgemeinste Form einer solchen Projektion lautet
\begin{align}
K: (\xi, \Dt Y^\mathfrak{g}|_{g}\cdot\xi)
	&\mapsto 
	(0, 
	\Dt Y^\mathfrak{g}|_{g}\cdot
	\Tt R_g|_e\cdot\xi
	+ \kappa(g, Y^\mathfrak{g}|_{g})\cdot\xi
	)
	\ .
\end{align}
Der \mydef{Konnektor} $K$ mit seiner durch $\kappa(g, Y^\mathfrak{g}|_{g}): \mathfrak{g}\rightarrow\mathfrak{g}$ bestimmten Projektion definiert einen Zusammenhang.

Hängt $\kappa$ für alle $g\in G$ linear von $Y^\mathfrak{g}|_{g}$ ab,
ist also
\begin{align}
	\kappa(g, Y^\mathfrak{g}|_{g})\cdot\xi
	&= \kappa|_g(Y^\mathfrak{g}|_{g}\otimes\xi)
\end{align}
bilinear in $Y^\mathfrak{g}|_{g}$ und $\xi$,
so definiert $\kappa$ einen \mydef{affinen Zusammenhang} auf $\Tt G$.
Die Identifizierung $\mathfrak{g}$ und $\Tt_0\mathfrak{g}$ liefert die zugehörige \mydef{kovariante Ableitung}
\begin{empheq}[box=\definitionbox]{align}	
	\nabla^\kappa_{\xi_g} Y|_g
	&=
	\Psi_1\big(
		g\,,\  
		\Dt Y^\mathfrak{g}|_{g}\cdot\xi_g
		+ \kappa|_g(Y^\mathfrak{g}|_{g}\otimes\Tt R_{g^{-1}}|_g\cdot\xi_g)
	\big)
	\nonumber\\
	&=
	\Tt R_g|_e\cdot
	\big(
		\Dt Y^\mathfrak{g}|_{g}\cdot\xi_g
		+ \kappa|_g(Y^\mathfrak{g}|_{g}\otimes\Tt R_{g^{-1}}|_g\cdot\xi_g)
	\big)
\end{empheq}
in Richtung des Tangentialvektors $\xi_g\in\Tt_gG$.
Zum auf diese Weise definierten Begriff der kovarianten Ableitung notieren wie einige Eigenschaften.

\paragraph{Christoffeltensor.}
Differenzen von linearen Zusammenhängen sind Tensoren.
Da bereits $\kappa=0$ einen Zusammenhang definiert, den $\nabla^+$-Zusammenhang, kann jeder linear Zusammenhang durch Vorgabe eines $(1,2)$-Christoffel\emph{tensors} $\kappa$ relativ zu $\nabla^+$ ausgedrückt werden.

\paragraph{Linearität.}
Der Zusammenhang $\nabla^\kappa_{\xi_g} Y|_g$ ist offensichtlich $\mathbb{R}$-linear bezüglich des Vektorfelds $Y$.

Sei $X$ ein weiteres Vektorfeld auf $G$.
Durch punktweise Auswertung von $X|_g$ wird das Vektorfeld
\begin{align*}
	\nabla^\kappa_{X} Y|_g
\end{align*}
gewonnen.
Die Linearität von $\nabla^\kappa$ bezüglich $\xi_g$ in jedem Tangentialraum impliziert, dass die Abhängigkeit von $X$, d.~h. die Richtungsabhängigkeit, $C^\infty$-linear ist.

\paragraph{Leipnizregel.}
Sei $f:G\rightarrow\mathbb{R}$ eine glatte Funktion.
Für $\xi_g\in\Tt_gG$ gilt
\begin{align*}
	\Dt(f\,Y^\mathfrak{g})|_g\cdot\xi_g
	&= 
	(\dt f|_g\cdot\xi_g)\,Y^\mathfrak{g}|_g
	+  f(g)\,\Dt Y^\mathfrak{g}|_g\cdot\xi_g
	\ .
\end{align*}
Durch Addition des bezüglich $Y^\mathfrak{g}$ $C^\infty$-linearen $\kappa$-Terms verifizieren wir mit $Y_g=\Tt R_g|_e\cdot Y^\mathfrak{g}|_g$ die Leipnizregel
\begin{empheq}[box=\resultbox]{align}	
	\nabla^\kappa_{\xi_g} (fY)|_g
	&= (\dt f|_g\cdot\xi_g)\,Y|_g
	+ f(g) \nabla^\kappa_{\xi_g} Y)|_g
	\ .
\end{empheq}

\subsection{$\nabla^+$-Zusammenhang}

Durch Wahl des Konnektors mit Christoffeltensor $\kappa=0$ wird der Zusammenhang $\nabla^+$ definiert.

\subsubsection{Definition des $\nabla^+$-Zusammenhangs}

Sei $G$ eine Lie-Gruppe und seien $X$ und $Y$ zwei Vektorfelder auf $G$ mit den Trivialisierungen~$X^\mathfrak{g}$ und $Y^\mathfrak{g}$.
Für $g\in G$ definieren wir den $\mathfrak{g}$-wertigen Ausdruck
\begin{empheq}[box=\definitionbox]{align}	
	\nabla^{\mathfrak{g},+}_{X^\mathfrak{g}}\,Y^\mathfrak{g}\big|_g
	&= \Dt Y^\mathfrak{g}|_g
	\cdot \Tt R_g|_e 
	\cdot X^\mathfrak{g}|_g
	\ .
\end{empheq}
Durch anschließenden Rücktransport zu~$g$,
\begin{empheq}[box=\definitionbox]{align}
	\nabla^+_{X}\,Y
	\big|_g
	&= \Tt R_g|_e
	\cdot
	\nabla^{\mathfrak{g},+}_{X^\mathfrak{g}}\,Y^\mathfrak{g}\big|_g
	\ ,
\end{empheq}
wird ein Zusammenhang auf $G$ definiert.
Er entspricht dem Zusammenhang $\nabla^\kappa$ mit $\kappa=0$ und erfüllt damit alle Voraussetzungen eines affinen Zusammenhangs.

\subsubsection{Eigenschaften des $\nabla^+$-Zusammenhangs}

Der Zusammenhang $\nabla^+$ auf $\Tt G$ mit seiner einfachen Struktur ergibt sich natürlich aus der Lie-Gruppenstruktur von $G$.

\paragraph{Ableitung rechtsinvarianter Vektorfelder.}
Für ein rechtsinvariantes Vektorfeld~$Y_\upsilon$ mit $Y_\upsilon|_e=\upsilon\in\mathfrak{g}$ ist $Y_\upsilon^\mathfrak{g}=\upsilon$ ein konstantes Feld $G\rightarrow\mathfrak{g}$, so dass die Ableitung $\Dt Y_\upsilon^\mathfrak{g}=0$ ist.
Es folgt
\begin{empheq}[box=\resultbox]{align}
\label{eq:nablaplus_rightinv}
	\nabla^+_{X}\,Y_\upsilon &= 0
	\ .
\end{empheq}

\paragraph{Vektorfelder bezüglich rechtsinvarianter Basis.}
Sei $\xi_1,\ldots,\xi_n$ eine Basis von $\mathfrak{g}=\Tt_eG$.
Bezüglich dieser Basis schreiben wir
\begin{align*}
	Y^\mathfrak{g}|_g
	&= \sum_{i=1}^n y_i(g)\,\xi_i
\end{align*}
mit glatten Koeffizientenfunktionen $y_i:G\rightarrow\mathbb{R}$.

Mit der Leipnizregel und, da Ableitungen rechtsinvarianter Felder verschwinden, folgt
\begin{empheq}[box=\resultbox]{align}
	\nabla^{\mathfrak{g},+}_{X^\mathfrak{g}}\,Y^\mathfrak{g}\big|_g
	&=
	\sum_{i=1}^n \dt y_i|_g\cdot X|_g
	\ .
\end{empheq}
Der Zusammenhang $\nabla^{+}$ reduziert sich also auf die Richtungsableitungen bezüglich der Komponentenfunktionen~$y_i$.

\paragraph{Torsion rechtsinvarianter Vektorfelder.}
Die Torsion 
\begin{align*}
	T^+(X,Y)
	&= \nabla^+_XY - \nabla^+_YX - [X,Y]
\end{align*}
des Zusammenhangs $\nabla^+$ berechnen wir zunächst für rechtsinvariante Vektorfelder.

Sind $X_\xi$, $Y_\upsilon$ Vektorfelder mit Werten $X_\xi|_e=\xi$ und $Y_\upsilon|_e=\upsilon$ im neutralen Element, so folgt aus Gl.~(\ref{eq:nablaplus_rightinv}) sofort, dass
\begin{empheq}[box=\resultbox]{align}
	T^+(X_\xi, Y_\upsilon)
	&=
	-[X_\xi, Y_\upsilon]
	= -[\xi, \upsilon]_\mathfrak{g}
	\ .
\end{empheq}

\paragraph{Torsion beliebiger Vektorfelder.}

\paragraph{Krümmung.}

\subsubsection{Paralleltransport und Geodätische des $\nabla^+$-Zusammenhangs}

\subsection{Lie-Gruppen als riemannsche Mannigfaltigkeiten}

\subsection{Liesche und riemannsche Exponentialabbildungen}

\section{Anwendung in der Navigation}




\begin{thebibliography}{9}

\bibitem{Gallot2005}
Sylvestre Gallot, Dominique Hulin, Jacques Lafontaine:
\textit{Riemannian Geometry}.
3rd Edition, Universitext, Springer, Berlin/Heidelberg, 2005.

\bibitem{StorchWiebe2005}
Uwe Storch, Hartmut Wiebe:
\textit{Lehrbuch der Mathematik, Band 4: Analysis auf Mannigfaltigkeiten -- Funktionentheorie -- Funktionalanalysis}.
Spektrum Akademischer Verlag, Heidelberg, 2005.

\end{thebibliography}



\end{document}